\section{Background}
\label{sec:preliminaries}

We review some preliminary facts on Coxeter groups and Hecke algebras relevant to the paper,
including the definition of Kazhdan--Lusztig cells, the definition of the $\la$-function, and the
classification of $\la(2)$-finite Coxeter systems. We then recall the notions of fully commutative
elements and generalized star operations, both of which will be essential to the study of
Kazhdan--Lusztig cells of $\la$-value 2.  

\subsection{Coxeter systems}
\label{sec:CoxeterSystems}
Throughout the paper, $(W,S)$ stands for a Coxeter system with a finite generating set $S$, Coxeter
group $W$ and Coxeter matrix $M=[m(s,t)]_{s,t\in S}$. Thus, we have $m(s,s)=1$ for all $s\in S$,
$m(s,t)=m(t,s)\in \Z_{\ge 2}\cup\{\infty\}$ for all distinct generators $s,t\in S$, and $W$ is the
group generated by $S$ subject to the relations $\{(ss')^{m(s,s')}=1, s,s'\in S, m(s,s')<\infty\}$.
It is well known that $st$ has order $m(s,t)$ for all distinct $s,t\in S$; in particular, $s$ and
$t$ commute if and only if~$m(s,t)=2$. 

The \emph{Coxeter diagram} of $(W,S)$ is the undirected diagram $G$ on vertex set $S$ where two
vertices $s,t$ are connected by an edge $\{s,t\}$ if and only if $m(s,t)\ge 3$. It follows that two
distinct generators in $S$ commute in $W$ if and only if they are not adjacent in $G$. For each edge
$\{s,t\}$ in $G$, we think of $m(s,t)$ as the \emph{weight} of the edge, call the edge \emph{simple}
if $m(s,t)=3$, and call the edge \emph{heavy} otherwise.  When drawing $G$, we label all heavy edges
by their weights, so that the defining data of $(W,S)$ can be fully recovered from the drawing.  The
system $(W,S)$ is \emph{irreducible} if the Coxeter diagram~$G$ is connected and \emph{reducible}
otherwise. 

\begin{remark}
\label{rmk:irreducible}
To simplify statements, we shall assume all Coxeter systems to be irreducible for the rest of the
paper.  The assumption is well justified as far as the $\la$-function is concerned, because the
$\la$-function behaves additively across connected components of the Coxeter graph. As a
consequence, Kazhdan--Lusztig cells of a particular $\la$-value in a reducible Coxeter system can be
easily obtained via cells of the same or lower $\la$-values in suitable irreducible Coxeter systems.
Since we are only interested in cells of $\la$-value 2, and cells of $\la$-values 1 or 0 in
irreducible Coxeter systems are well understood as explained in the introduction, it suffices to
consider only irreducible Coxeter systems as we study cells of $\la$-value 2 in this paper; see also
\cite[Theorem 1.3]{GreenXu}.
\end{remark}

For every subset $J$ of $S$, the subgroup $W_J:=\ip{s:s\in J}$ of $W$ is called the \emph{parabolic
subgroup generated by $J$}. Note that if for another Coxeter system $(W',S')$ there is a bijection
$f: S'\ra J$ with $m(s,t)=m(f(s),f(t))$ for all $s,t\in S'$ (in other words, if the Coxeter diagram
of $(W',S')$ is isomorphic to the subgraph of the Coxeter diagram of $(W,S)$ induced by $J$), then
this bijection naturally extends to a group isomorphism from $W'$ to $W_J$. For example, in the
notation of Proposition~\ref{prop:facts01}.(4) the Coxeter group $A_{n-1}$ naturally embeds into $B_{n}$ for
all $n\ge 2$ in the sense that $A_{n-1}$ is isomorphic to the parabolic subgroup generated by the
set $J=\{2,\dots,n\}$. 

Let $S^*$ be the free monoid on $S$, viewed as the set of words on the alphabet $S$. To distinguish
words in $S^*$ from elements of $W$, we denote each word with an underline.  For example, the words
$\ul u=st, \ul v=ts\in S^*$ represent the same element $w=st=ts$ in~$W$ if $s,t$ are commuting
distinct generators in $S$.  For each $w\in W$, the words $\ul w\in S^*$ that express $w$ with a
minimum number of letters are called the \emph{reduced words} of~$w$; that minimum number is called
the \emph{length} of $w$ and written $l(w)$. More generally, we call a factorization $w=w_1w_2\dots
w_k$ in $w$ \emph{reduced} if $l(w)=\sum_{i=1}^k l(w_i)$. In this paper, the notation $w=w_1\cdot
w_2\cdot \dots \cdot w_k$ will always indicate that the factorization $w=w_1w_2\dots w_k$ is
reduced.  Whenever we have $z=x\cdot y$ for elements $x,y,z\in W$, we write $x\le^R z$ and $y\le^L
z$. The relations $\le^R$ and $\le^L$ define two partial orders called the \emph{right weak Bruhat
order} and \emph{left weak Bruhat order} on $W$, respectively.

For each $w\in W$, we define a \emph{left descent} of $w$ to be a generator $s\in S$ such that
$l(sw)<l(w)$, and we denote the set of left descents of $w$ by $\call(w)$. Similarly, a \emph{right
descent} of $w$ is a generator $s\in S$ such that $l(ws)<l(w)$, and we denote the set of right
descents of $w$ by $\calr(w)$. Descents can be described via reduced words: it is well known that
for all $s\in S$ and $w\in W$, we have $s\in \call(w)$ if and only if some reduced word of $w$
begins with the letter $s$, and $s\in \calr(w)$ if and only if some reduced word of~$w$ ends with
the letter $s$.  

We end the subsection by recalling the Matsumoto--Tits Theorem. Define a \emph{braid relation} to be
a relation of the form $sts\cdots=tst\cdots$ where $s,t\in S, 2\le m(s,t)<\infty$, and both sides
contain $m(s,t)$ factors. Then the theorem states that for all $w\in W$, every pair of reduced words
of $w$ can be obtained from each other by applying a finite sequence of braid relations. It follows
that every two reduced words of an element~$w$ contain the same set of letters. We call the set the
\emph{support} of $w$ and denote it by~$\supp(w)$.

\subsection{Cells and the \texorpdfstring{$\la$}{a}-function}
\label{sec:HeckeAlgebras}
Let $\cala=\Z[v,v\inverse]$ and let $H$ be the \emph{Hecke algebra} of a Coxeter system $(W,S)$.  We
recall that $H$ is a unital, associative $\cala$-algebra given by a certain presentation and that
$H$ has an $\cala$-linear basis $\{C_w:w\in W\}$ called the \emph{Kazhdan--Lusztig} basis; see
\cite[\S 2.2]{XuSubregular}. The Kazhdan--Lusztig basis gives rise to Kazhdan--Lusztig cells and the
$\la$-function as follows.


Let $x,y\in W$. For all $z\in W$, let $D_z: H\ra\cala$ be the unique linear map such that
$D_z(C_w)=\delta_{z,w}$ for all $w\in W$, where $\delta$ is the Kronecker delta symbol.  Write~$x\prec_L y$ 
or $x\prec_R y$ if $D_x(C_sC_y)\neq 0$ or $D_{x}(C_yC_s)\neq 0$ for some $s\in S$,
respectively, and write $x\prec_{LR}y$ if either $x\prec_L y$ or $x\prec_R y$.  The preorders
$\le_L,$ $\le_R,$ $\le_{LR}$ defined as the reflexive, transitive closures of the relations
$\prec_L,\prec_R, \prec_{LR}$ naturally generate equivalence relations $\sim_L, \sim_R, \sim_{LR}$,
and we call the resulting equivalence classes the  \emph{left, right and two-sided Kazhdan--Lusztig
cells} of $W$, respectively.  For example, we have that~$x\sim_L y$, i.e. $x$ and $y$ are in the same
left Kazhdan--Lusztig cell, if and only if $x\le_L y$ and $y\le_L x$, which happens if and only if
there are sequences $x=w_0,w_1,\dots, w_n=y$ and $y=z_0,z_1,\dots, z_k=x$ where $n,k\ge 0$ and
$w_i\prec_L w_{i+1}$, $z_j\prec_L z_{j+1}$ for all $0\le i<n, 0\le j<k$.  For brevity, from now on we
will refer to Kazhdan--Lusztig cells simply as ``cells'', and we call both left and right cells
\emph{one-sided} cells. By the definition of $\prec_{LR}$, each two-sided cell of $W$ is a union of
left cells, as well as a union of right cells. 

To define the $\la$-function from the basis $\{C_w:w\in W\}$, consider the structure constants
$h_{x,y,z}\in \cala \, (x,y,z\in W)$ such that 
\begin{equation}
    \label{eq:h} C_x C_y =\sum_{z\in W} h_{x,y,z} C_z 
\end{equation} for all $x,y\in W$. By \S 15.2 of \cite{LusztigHecke}, for each $z\in W$ there exists
    a unique integer $\la(z)\ge 0$ such that 
\begin{enumerate}
    \item[(a)] $h_{x,y,z}\in v^{\la(z)}\Z[v\inverse]$ for all $x,y\in W$;
    \item[(b)] $h_{x,y,z}\notin v^{\la(z)-1}\Z[v\inverse]$ for some $x,y\in W$. 
\end{enumerate} The assignment $z\mapsto \la(z)$ defines the function $\la: W\ra \Z_{\ge 0}$.

In the sequel, we will make extensive use of the following standard results on cells and the
$\la$-function. We note that the proofs of Parts (5)--(9) of the proposition either directly or
indirectly rely on the use of Lusztig's boundedness conjecture \cite[\S 13.4]{LusztigHecke}. The
conjecture states that for the standard basis $\{T_w:w\in W\}$ of the Hecke algebra (see \cite[\S
2.2]{XuSubregular}) and the structure constants $f_{x,y,z}$ such that $T_xT_y=\sum_{z\in
W}f_{x,y,z}T_z$ for all $x,y\in W$, there is a nonnegative integer $N$ such that $v^{-N}f_{x,y,z}$
is an element of the ring $\Z[v\inverse]$ for all $x,y,z\in W$.  In the following proposition, Parts
(5)--(9) contain statements from conjectures P9--P12 and P14 in \cite[\S 14]{LusztigHecke}, and the
proofs of these five conjectures in turn rely on Lusztig's boundedness conjecture (despite the fact
    that conjectures P1--P6 from \cite[\S 14]{LusztigHecke} are known to hold in the setting of our
current paper, where the Hecke algebra has ``equal parameters'', without reliance on the boundedness
conjecture); see also \cite[\S 15.1]{LusztigHecke}, \cite[\S 14.3]{BonnafeCells} and \cite[Theorem
15.2.5]{BonnafeCells}.

\begin{prop}
     \label{prop:Facts}
     Let $(W,S)$ be a Coxeter system and let $x,y\in W$.
     \begin{enumerate}
         \item If $x\le^L y$ then $y\le_L x$; if $x\le^R y $ then $y\le_R x$.
         \item If $x$ is a product of $k$ pairwise commuting generators in $S$, then $\la(x)=k$.
         \item We have $\call(x)=\call(y)$ if $x\sim_R y$, and $\calr(x)=\calr(y)$ if $x\sim_L y$;
         \item For each left cell $L$ in $W$, the set $L\inverse=\{z\inverse: z\in L\}$ is a right cell; for each
             right cell $R$ in $W$, the set $R\inverse=\{z\inverse:z\in R\}$ is a left cell. 
         \item If $x\le_{LR} y$, then we have $\la(x)\ge \la(y)$. In particular, if $x\sim_{LR}y$, then
             we have $\la(x)=\la(y)$. Moreover, if $x\le_L y$ but $x\not\sim_L y$, or if $x\le_R y$ but
             $x\not\sim_R y$, then we have $\la(x)>\la(y)$.
         \item We have $z\sim_{LR} z\inverse$ for all $z\in W$.
         \item Let $J\se S$ and consider the corresponding parabolic subgroup $W_J$ of $W$.
             Let $\la_J$ denote the $\la$-function associated to the Coxeter system $(W_J,J)$. If 
             $x\in W_J$, then
             the $\la$-value $\la_J(x)$ computed in the system $(W_J, J)$ equals the value
             $\la$-value $\la(x)$ computed in the system $(W,S)$.
         \item Let $L_1,L_2$ be two left cells in the same two-sided
             cell. Then the intersection $L_1\cap (L_2\inverse)$ is
             nonempty.
         \item If $x\sim_{LR}y$, then $W$ contains an element $z$ such that $z\sim_L x$ and $z\sim_R y$. 
     \end{enumerate}
        \end{prop}
        \begin{proof}
Parts (1)--(7) are due to Lusztig and can be found in \cite{LusztigHecke}: Parts (1)--(2) follow
immediately from Theorem 6.6 and Proposition 2.6; Parts (3)--(4) are proved in Section 8; the
statements in Parts (5)--(7) appear in Sections 14 as conjectures P9--P12 and P14.  Part~(8) is
Lemma 2.4 of \cite{BO}. (The lemma relies on the construction of the so-called $J$-ring of $(W,S)$
from \cite[\S 18]{LusztigHecke}, which in turn depends on Lusztig's boundedness conjecture.)  To see
Part~(9), note that if $x\sim_{LR} y$ then $x\sim_{LR} y\inverse$ by Part~(6); therefore the left
cell of $x$ intersects the inverse of the left cell of $y\inverse$ nontrivially by Part~(8).  This
inverse is precisely the right cell of $y$ by Part~(4), so any element $z$ in the intersection
satisfies $z\sim_L x$ and $z\sim_R y$.  
    \end{proof}

\begin{definition}
    For each Coxeter system $(W,S)$ and each integer $n\ge 0$, we let $W_n=\{w\in W:\la(w)=n\}$, and we say that
    $(W,S)$ is \emph{$\la(n)$-finite} if $W_n$ is finite. 
    \label{def:an-finite}
\end{definition}
 
By Proposition~\ref{prop:Facts}.(5), for each integer $n\ge 0$ the set $W_n$ defined above is a union of
two-sided cells, of left cells, and of right cells.  We are interested in the cells in $W_2$ within
$\la(2)$-finite Coxeter systems.  Our motivation partly comes from the following facts about $W_0$
and $W_1$. 

\begin{prop}
    \label{prop:facts01}
    Let $(W,S)$ be a Coxeter system with Coxeter diagram $G$. \tul Recall from Remark~\ref{rmk:irreducible} that
    we assume $(W,S)$ is irreducible.\tur
    \begin{enumerate}
        \item We have $W_0=\{1\}$. In particular, $W_0$ is finite and simultaneously
            the unique left, right and two-sided cell of $\la$-value 0. 
        \item We have  $W_1=\{w\in W:w\neq 1 \text{ and $w$ has a unique reduced word}\}$. The
            set $W_1$ is finite \textup{(}i.e. the
                system $(W,S)$
            is $\la(1)$-finite\textup{)} if and only if $G$ is acyclic, contains no edge with infinite
            weight, and contains at most one heavy edge with finite weight.  
        \item The set $W_1$ is itself a two-sided cell and hence the unique two-sided cell of $\la$-value 1 in
            $W_1$. The left cells in $W_1$ are in bijection with $S$ and given by the sets \[ L_{s}:=\{z\in
            W_1:s\le^L z\} \] while the right cells in $W_1$ are in bijection with $S$ and given by the sets \[
            R_{s}:=\{z\in W_1:s\le^R z\} \] where $s\in S$. 
        \item If $G$ contains a cycle, then $W$ is $\la(2)$-finite if and only if $G$ is complete. If~$G$ is
            acyclic, then $W$ is $\la(2)$-finite if and only if $G$ is one of the graphs in Figure~\ref{fig:finite
            E}, where $n$ equals the number of vertices in the Coxeter diagram and there are $q$ and $r$ vertices
            strictly to the left and the right of the trivalent vertex in $E_{q,r}$. 
            \begin{figure}
                \begin{tikzpicture}

                    \node(1) {$A_n$};
                    \node[main node] (11) [right=0.5cm of 1] {};
                    \node[main node] (12) [right=1cm of 11] {};
                    \node[main node] (13) [right=1cm of 12] {};
                    \node[main node] (15) [right=1.5cm of 13] {};
                    \node[main node] (16) [right=1cm of 15] {};
                    \node (17) [right=0.5cm of 16] {$(n\ge 1)$}; 

                    \node (a1) [below =0.1cm of 11] {\tiny{$1$}};
                    \node (a2) [below =0.1cm of 12] {\tiny{$2$}};
                    \node (a3) [below =0.1cm of 13] {\tiny{$3$}};
                    \node (a4) [below =0.05cm of 15] {\tiny{$(n-1)$}};
                    \node (a5) [below =0.1cm of 16] {\tiny{$n$}};

                    \path[draw]
                    (11)--(12)--(13)
                    (15)--(16);

                    \path[draw,dashed]
                    (13)--(15);

                    \node(2) [below=1cm of 1] {$B_n$};

                    \node[main node] (21) [right=0.5cm of 2] {};
                    \node[main node] (22) [right=1cm of 21] {};
                    \node[main node] (23) [right=1cm of 22] {};
                    \node[main node] (25) [right=1.5cm of 23] {};
                    \node[main node] (26) [right=1cm of 25] {};
                    \node (27) [right=0.5cm of 26] {$(n\ge 2)$}; 

                    \node (b1) [below =0.1cm of 21] {\tiny{$1$}};
                    \node (b2) [below =0.1cm of 22] {\tiny{$2$}};
                    \node (b3) [below =0.1cm of 23] {\tiny{$3$}};
                    \node (b4) [below =0.05cm of 25] {\tiny{$(n-1)$}};
                    \node (b5) [below =0.1cm of 26] {\tiny{$n$}};

                    \path[draw]
                    (21) edge node [above] {\tiny{$4$}} (22)
                    (22)--(23)
                    (25)--(26);

                    \path[draw,dashed]
                    (23)--(25);

                    \node (c) [below=1cm of 2] {$\tilde C_{n-1}$};
                    \node[main node] (c1) [right=0.3cm of c] {};
                    \node[main node] (c2) [right=1cm of c1] {};
                    \node[main node] (c3) [right=1cm of c2] {};
                    \node[main node] (c5) [right=1.5cm of c3] {};
                    \node[main node] (c6) [right=1cm of c5] {};
                    \node (c7) [right=0.5cm of c6] {$(n\ge 5)$}; 

                    \node (b11) [below =0.1cm of c1] {\tiny{$1$}};
                    \node (b22) [below =0.1cm of c2] {\tiny{$2$}};
                    \node (b33) [below =0.1cm of c3] {\tiny{$3$}};
                    \node (b44) [below =0.05cm of c5] {\tiny{$(n-1)$}};
                    \node (b55) [below =0.1cm of c6] {\tiny{$n$}};

                    \path[draw]
                    (c1) edge node [above] {\tiny{$4$}} (c2)
                    (c2)--(c3)
                    (c5) edge node [above] {\tiny{$4$}} (c6);

                    \path[draw,dashed]
                    (c3)--(c5);


                    \node (3) [below=1.5cm of c] {$E_{q,r}$};

                    \node[main node] (31) [right=0.4cm of 3] {};
                    \node[main node] (32) [right=1cm of 31] {};
                    \node[main node] (34) [right=1.5cm of 32] {};
                    \node[main node] (38) [right=1.5cm of 34] {};
                    \node[main node] (39) [right=1cm of 38] {};
                    \node (40) [right=0.5cm of 39] {$(r\ge q\ge 1)$}; 
                    \node[main node] (a) [above=0.8cm of 34] {};

                    \path[draw]
                    (34)--(a)
                    (31)--(32)
                    (38)--(39);
                    \path[draw,dashed]
                    (32)--(34)--(38);

                    \node (c11) [below =0.1cm of 34] {\tiny{$0$}};
                    \node (c22) [below =0.1cm of 31] {\tiny{$-q$}};
                    \node (c33) [below =0.05cm of 32] {\tiny{$-(q-1)$}};
                    \node (c44) [below =0.05cm of 38] {\tiny{$(r-1)$}};
                    \node (c55) [below =0.1cm of 39] {\tiny{$r$}};
                    \node (c66) [right =0.02cm of a] {\tiny{$v$}};


                    \node(4) [below=1cm of 3] {${F}_n$};

                    \node[main node] (41) [right=0.5cm of 4] {};
                    \node[main node] (42) [right=1cm of 41] {};
                    \node[main node] (43) [right=1cm of 42] {};
                    \node[main node] (45) [right=1.5cm of 43] {};
                    \node[main node] (46) [right=1cm of 45] {};
                    \node (47) [right=0.5cm of 46] {$(n\ge 4)$}; 

                    \node (d11) [below =0.1cm of 41] {\tiny{$1$}};
                    \node (d22) [below =0.1cm of 42] {\tiny{$2$}};
                    \node (d33) [below =0.1cm of 43] {\tiny{$3$}};
                    \node (d44) [below =0.05cm of 45] {\tiny{$(n-1)$}};
                    \node (d55) [below =0.1cm of 46] {\tiny{$n$}};


                    \path[draw]
                    (41)--(42)
                    (42) edge node [above] {\tiny{$4$}} (43)
                    (45)--(46);
                    \path[draw,dashed]
                    (43)--(45);

                    \node(h) [below=1cm of 4] {$H_n$};

                    \node[main node] (h1) [right=0.45cm of h] {};
                    \node[main node] (h2) [right=1cm of h1] {};
                    \node[main node] (h3) [right=1cm of h2] {};
                    \node[main node] (h5) [right=1.5cm of h3] {};
                    \node[main node] (h6) [right=1cm of h5] {};
                    \node (h7) [right=0.5cm of h6] {$(n\ge 3)$}; 

                    \node (bb1) [below =0.1cm of h1] {\tiny{$1$}};
                    \node (bb2) [below =0.1cm of h2] {\tiny{$2$}};
                    \node (bb3) [below =0.1cm of h3] {\tiny{$3$}};
                    \node (bb4) [below =0.05cm of h5] {\tiny{$(n-1)$}};
                    \node (bb5) [below =0.1cm of h6] {\tiny{$n$}};

                    \path[draw]
                    (h1) edge node [above] {\tiny{$5$}} (h2)
                    (h2)--(h3)
                    (h5)--(h6);
                    \path[draw,dashed]
                    (h3)--(h5);

                    \node(i) [below=1cm of h] {$I_2(m)$};

                    \node[main node] (i1) [right=0.25cm of i] {};
                    \node[main node] (i2) [right=1cm of i1] {};
                    \node (i3) [right=0.5cm of i2] {$(5\le m\le \infty)$};


                    \node (bb1) [below =0.1cm of i1] {\tiny{$1$}};
                    \node (bb2) [below =0.1cm of i2] {\tiny{$2$}};

                    \path[draw]
                    (i1) edge node [above] {\tiny{$m$}} (i2);
                \end{tikzpicture}
                \caption{Irreducible $\la(2)$-finite Coxeter groups with acyclic diagrams.}
                \label{fig:finite E}
            \end{figure}   
    \end{enumerate}
\end{prop}
\begin{proof} Part~(1) is proved in \cite[\S 13]{LusztigHecke}.  Let 
\[
C= \{w\in W:w\neq 1 \text{ and
      $w$ has a unique reduced word}\}.\] Then the first assertion in Part~(2), the assertion that
      $W_1=C$, is proved in \cite[Corollary 3.1]{XuSubregular}. Given the equality $W_1=C$, the
      second assertion in Part~(2) follows immediately from \cite[Proposition
      3.8(h)]{LusztigSubregular}.  Part~(3) follows from  \cite[Proposition~3.8(c)]{LusztigSubregular} 
      and Proposition~\ref{prop:Facts}(4).  Part~(4) is the main result of
      \cite{GreenXu}. 
\end{proof}

\begin{remark}[Labeling conventions]
    \label{rmk:labels}
    For each $\la(2)$-finite Coxeter system $(W,S)$, we will adhere to the labeling of the set $S$
    from Figure~\ref{fig:finite E} for the rest of the paper. For example, the generators in $B_n$ will
    always be labeled $1,2,\dots,n$, with $\{1,2\}$ forming the unique heavy edge, and generators in
    $E_{q,r}$ will always be labeled by the numbers $-q,-q+1,\dots, 0,\dots, r$ and the letter $v$
    under the assumption that $r\ge q$, which causes no loss of generality. Note that in the cases
    where $q=1$, $(q,r)=(2,2)$, $(q,r)=(2,3)$, $(q,r)=(2,4)$, $(q,r)=(3,3)$ and $(q,r)=(2,5)$, the Coxeter
    group of type $E_{q,r}$ coincides with the Weyl or affine Weyl group of type $D_{r+3}, E_6, E_7,
    E_8, \tilde E_7$ and $\tilde E_8$, respectively; the group $F_5$ coincides with the affine Weyl
    group of type $\tilde F_4$.
\end{remark}

As we study cells within the set $W_2$ for $\la(2)$-finite Coxeter groups, Part~(4) of
Proposition~\ref{prop:Facts} allows us to consider only right cells because each left cell $L$ is the setwise
inverse of a right cell $R:=L^{-1}$.  Furthermore, Part~(6) of the same proposition implies that the
set $R$ is not only a right cell but also a right cell in the same two-sided cell as $L$. The
following corollary, which we will use to count cells in $W_2$, is now immediate. 

\begin{corollary}
\label{coro:intersection} Let $E$ be a two-sided cell in a Coxeter group $W$. Let $\call\calc$ and
      $\calr\calc$ be the collections of left cells and right cells in $E$, respectively. Then
      $\call\calc=\{R\inverse: R\in \calr\calc\}$. In particular, for each $R\in \calr\calc$ we have
      \[ R=\bigsqcup_{L\in\call\calc}R\cap L=\bigsqcup_{R'\in \calr\calc}R\cap (R')\inverse \] and
      hence \[ \abs{R}=\sum_{L\in\call\calc}\Abs{R\cap L}=\sum_{R'\in \calr\calc}\Abs{R\cap
      (R')\inverse}.  \]
\end{corollary}

The above corollary suggests that intersections of left and right cells in a two-sided cell are
important. This motivates the following definition.
\begin{definition} 
    We define a \emph{zero-sided cell} in a Coxeter system $(W,S)$ to be a set of the form $R\cap L$
    where $R$ is a right cell and $L$ is a left cell in the same two-sided cell as $R$.  For
    brevity, we also refer to two-sided, one-sided and zero-sided cells of $W$ as \emph{2-cells},
    \emph{1-cells} and \emph{0-cells}.
    \label{def:0cells}
\end{definition}

\begin{example}
    \label{eg:b4}
    Let $(W,S)$ be the Coxeter system of type $B_4$, which is $\la(2)$-finite by
    Proposition~\ref{prop:Facts}.(10). Later we will see that $W_2$ is the unique two-sided cell of $\la$-value
    2 in $W$; moreover, the cell $W_2$ can be partitioned into 6 right cells as well as into 6 left
    cells. Under a certain labeling of the right and left cells respectively as $R(i)$ and $L(j)$
    for $1\le i,j\le 6$, the sizes of the 0-cells in $W_2$ are given by Table~\ref{tab:b4example}, where
    the entry in the $R(i)$-row, $L(j)$-column equals the size of the 0-cell $R(i)\cap L(j)$. By
    Corollary~\ref{coro:intersection}, we may count each 1-cell $R(i)$ by summing the entries in the
    corresponding row, and the we may count the 2-cell $W_2$ by summing all entries in the table.
    It follows that $\abs{R(i)}=10$ if $1\le i\le 4$, that $\abs{R(i)}=8$ if $5\le i\le 6$, and that
    $\abs{W_2}=56$.
    \begin{table} \centering
        \begin{tabular}{|c|c|c|c|c|c|c|} \hline & $L(1)$ & $L(2)$ & $L(3)$ & $L(4)$ & $L(5)$ &
              $L(6)$\\\hline $R(1)$ & 2& 2&  2 & 2 &  1&  1\\\hline $R(2)$ & 2& 2&  2 & 2 &  1&
              1\\\hline $R(3)$ & 2& 2&  2 & 2 & 1&  1\\\hline $R(4)$ & 2& 2&  2 & 2 &  1&  1\\\hline
              $R(5)$ & 1& 1&  1 & 1 &  2&  2\\\hline $R(6)$ & 1& 1&  1 & 1 &  2&  2\\\hline
        \end{tabular} \vspace{1em} \caption{Sizes of 0-cells of $\la$-value 2 in type $B_4$}
        \label{tab:b4example}
    \end{table}
\end{example}

\begin{remark}[Symmetry]
    \label{rmk:symmetry}
    As is evident in Propositions~\ref{prop:Facts} and~\ref{prop:facts01}, there is a large amount
    of ``left-right symmetry'' in the properties of Coxeter group elements and Kazhdan--Lusztig
    cells.  The rest of the paper contains many more pairs of ``left-right symmetric'' notions and
    assertions, such as the two statements in Proposition~\ref{prop:StarOpAndCells}. For such pairs, we will
    often formulate one notion or assertion carefully and leave the formulation of the other to the
    reader (as we actually already did in Corollary~\ref{coro:intersection} and Example~\ref{eg:b4}), or prove one
    of the assertions and invoke symmetry as justification for the other.  The reader may assume
    that the obvious symmetry always works as expected. 
\end{remark}

\subsection{Fully commutative elements}
\label{sec:FC}

Recall that a braid relation in a Coxeter system is a relation of the form $st\dots =ts\dots$ where
$s,t\in S, 2\le m(s,t)<\infty$ and both sides contain $m(s,t)$ factors. We call the relation a
\emph{commutation relation} if $m(s,t)=2$ and call each side of the relation a \emph{long braid} if
$m(s,t)\ge 3$. An element $w\in W$ is called \emph{fully commutative}, or \emph{FC}, if every pair
of reduced words of $w$ can be connected by a finite sequence of commutation relations. Proposition
2.1 of \cite{FC} gives a well-known ``word criterion for full commutativity'': we have $w$ if FC if
and only if no reduced word of $w$ contains a long braid as a subword. Here and henceforth, by a
subword of a word $s_1\dots s_q$ we always mean a contiguous subword, i.e. a word of the form
$s_is_{i+1}\dots s_{j-1}s_{j}$ for some $1\le i\le j\le q$.

We denote the set of all FC elements in $W$ by $\fc(W)$. We shall study cells of $\la$-value 2 via
FC elements because of the following fact:

\begin{prop}[{\cite[Proposition 3.9]{GreenXu}}] \label{prop:a2impliesFC} 
    Let $w\in W$. If $\la(w)=2$, then $w$ is FC.  
\end{prop} 
 
We review two important tools for studying FC elements.  The first is a canonical form called the
\emph{Cartier--Foata form}. For an FC element $w$, this refers to the reduced word $w=w_p\cdot \dots
\cdot w_2\cdot w_1$ satisfying the following properties:
\begin{enumerate}
    \item  For all $1\le i\le p$, the element $w_i$ is a product of pairwise commuting generators in
         $S$;
    \item For all $1<i\le p$, every generator $t\in \supp(w_i)$ fails to commute with some generator
         $s\in \supp(w_{i-1})$.
\end{enumerate} 
The factors $w_1,w_2,w_3,\ldots$ can be obtained inductively as the product of the right descents of
the elements $x_1=w, x_2=x_1w_1, x_3=x_2w_2, \ldots $, respectively, and the factorization $w=w_p\dots w_1$
is unique up to the re-ordering of the generators appearing in~$w_i$ for each $i$; see
\cite{GreenStarReducible}.  We call each $w_i$ the \emph{$i$-th layer} of $w$. When the generators
in $S$ are labeled by integers, we shall insist that within each layer we order the generators in
increasing order from left to right. For example, in the group $W=A_5$ the element $w=1532\in
\fc(W)$ has Cartier--Foata form $w=w_2w_1$ where $w_1=25$ and $w_2=13$.  For more on the
Cartier--Foata form, see \cite{GreenStarReducible}.

Our second tool for studying an FC element $w$ is a labeled poset called a {heap}. To define it, we
first define the \emph{heap of a word} $\underline w=s_1\dots s_q\in S^*$ to be the labeled poset
$H(\underline w):=([q], \preceq)$ where the underlying set is $[q]=\{1,2,\dots,q\}$, the partial
order $\preceq$ is the reflexive, transitive closure of the relation $\prec$ defined by \[ i\prec j
\text{\quad if $i< j$ and $m(s_i,s_j)\neq 2$}, \] and the label of the element $i$ is $s_i$ for each
$i\in [q]$.  The heaps of two words related by a commutation relation are isomorphic as labeled
posets in the sense that there exists a poset isomorphism $f:H({\ul w})\ra H({\ul w'})$ such that
$f(i)$ and $i$ have the same label for all $i\in H({\ul w})$; see \cite[\S 2.2]{FC}.  Thus, it makes
sense to define the \emph{heap of an element} $w\in W$, up to isomorphism, to be the heap of any
reduced word of $w$; we denote this heap by $H(w)$. The following ``heap criterion for full
commutativity'' is a well-known analog of the word criterion mentioned earlier. 

\begin{prop}[{\cite[Proposition 3.3]{FC}}]
    \label{prop:FCCriterion}
    Let  $\underline w= s_1s_2\cdots s_q\in S^*$. Then $\underline w$ is the reduced word of a fully commutative
    element in $W$ if and only if the heap $H(\underline w)$ satisfies the following conditions: 
    \begin{enumerate}
        \item there is no covering relation $i\prec j$ such that $s_i=s_j$;
        \item there is no convex chain $i_1\prec i_2\prec \ldots \prec i_m$ in $H(\underline w)$ with
             $s_{i_1}=s_{i_3}=\cdots = s$ and $s_{i_2}=s_{i_4}=\cdots=t$ where $s,t\in S$ and $m = m(s,t)\ge 3$.
    \end{enumerate} 
\end{prop}

\begin{remark}
    Let $w\in \fc(W)$.  By the  definition of $H(w)$, the left and right descents of $w$ are exactly
    the labels of the minimal and maximal elements in $H(w)$, respectively. In particular, the
    support of the first layer in the Cartier--Foata form of $w$ coincides with the set of the
    labels of the maximal elements in $H(w)$, because they both equal~$\calr(w)$.
    \label{rmk:heapsAndWords}
\end{remark}

\begin{remark}
    \label{rmk:chainsAndWalks}
    By the definition of a heap, for every chain of coverings $i_1\prec i_2 \prec \dots \prec i_k$ in the
    heap of an FC element $w=s_1\dots s_q$, the generators $s_{i_j}$ and $s_{i_{j+1}}$  must be
    distinct and adjacent in the Coxeter diagram. In other words, the sequence of generators
    $s_{i_1},\dots, s_{i_k}$ forms a \emph{walk} from $s_{i_1}$ to $s_{i_k}$ on the Coxeter diagram
    in the usual graph theoretical sense. We also recall, for later use, that such a walk has
    \emph{length} $(k-1)$ and that a \emph{path} from a vertex $v$ to a vertex~$u$ on a graph is a
    walk from $v$ to~$u$ of minimal length. On an acyclic and connected graph, there is a unique
    path from $u$ to $v$ for any two vertices $u,v$.  
\end{remark}

There is an intuitive way to visualize the heap of a word $\underline w=s_1\dots s_q\in S^*$ in the
lattice $S\times \Z_{\ge 0}$. Here, we think of the elements of $S$ as \emph{columns} and the
elements of $\Z_{\ge 0}$ as \emph{levels}, and we say two columns $s,t$ \emph{commute} if they
commute as generators, i.e. if $m(s,t)= 2$. To embed $H(\underline w)$ into the lattice, we read
the letters in $\underline w$ from left to right, and for the $j$-th letter we drop a vertex $p_j$
in the column $s_j$ to the lowest level possible subject to the condition that $p_j$ should fall
above every vertex $p_i$ that has been dropped into a column that does not commute with $s_j$, i.e.
above every vertex $p_i$ for which $i\prec j$. In addition, if a column $s$ does not commute with
$s_j$ and contains at least one vertex $p_i$ with $i<j$, i.e. if $i\prec j$, then we draw an edge
to connect $p_j$ to the highest such vertex, i.e. to the vertex $p_i$ in column $s$ with $i$
maximal. It follows that the vertices $p_1,\dots, p_q$ and the edges of the form $\{p_i,p_j\}$ where
$i\prec j$ form exactly the Hasse diagram of the heap $H(\underline w)$. 

When drawing the embedding of $H(\underline w)$, it is customary to not draw the columns or levels
and to label each point $p_i$ simply by $s_i$. This ensures that the isomorphic heaps arising from
the reduced words of $w$ are given by the same graph, the Hasse diagram of the heap $H(w)$. For
example, in Figure~\ref{fig:HeapExample} the picture on the right shows the heap $H(w)$ of the
element $w = abcadb$ in the Coxeter group whose Coxeter diagram is drawn on the left. 
\begin{figure} \centering
    \begin{tikzpicture} \node[main node] (a) {}; \node[main node] (b) [right=1cm of a] {};
          \node[main node] (c) [right=1cm of b] {}; \node[main node] (d) [right=1cm of c] {};

        \node (aa) [below=0.2cm of a] {$a$};
        \node (bb) [below=0.1cm of b] {$b$};
        \node (cc) [below=0.2cm of c] {$c$};
        \node (dd) [below=0.1cm of d] {$d$};

        \path[draw]
        (b)--(c)--(d)
        (a) edge node [above] {$4$} (b);

        \node[main node] (1) [below right=2cm and 5cm of a] {};
        \node[main node] (3) [above right=1cm and 1cm of 1] {};
        \node[main node] (4) [above left=1cm and 1cm of 3] {};
        \node[main node] (5) [above right=1cm and 1cm of 3] {};
        \node[main node] (6) [above left=1cm and 1cm of 5] {};
        \node[main node] (7) [above right=1cm and 1cm of 5] {};

        \node (11) [above=0.1cm of 1] {$a$};
        \node (33) [above=0.1cm of 3] {$b$};
        \node (44) [above=0.1cm of 4] {$a$};
        \node (55) [above=0.1cm of 5] {$c$};
        \node (66) [above=0.1cm of 6] {$b$};
        \node (77) [above=0.1cm of 7] {$d$};

        \path[draw]
        (1)--(3)--(4)--(6)--(5)--(7)
        (3)--(5);
    \end{tikzpicture}
    \caption{Lattice embedding of a heap.}
    \label{fig:HeapExample}
\end{figure}

Heaps can help compute $\la$-values. Let $n(w)$ be the  \emph{width} of the poset $H(w)$, i.e. let 
\[
    n(w)= \max\{\abs{A}:A \text{\; is an antichain in $H(w)$}\}.
\]
Then $n(w)$ bounds and may equal $\la(w)$ (see also Remark~\ref{rmk:a.vs.n}):

\begin{prop}[{\cite{Shi}}]\label{prop:a=n}
    Let $w\in \fc(W)$. Then
    \begin{enumerate}
        \item we have $\la(w)\ge n(w)$;
        \item we have $\la(w)=n(w)$ if $W$ is a Weyl or affine Weyl group.
    \end{enumerate}
\end{prop}

\begin{proof}
    The equality in (2) is the main result, Theorem 3.1, of \cite{Shi}. The inequality in~(1) essentially follows
    the results in \S 1.3 of the same paper, but let us give a short proof in our notation: the definition of
    heaps implies that $w$ admits a reduced factorization of the form $w=x\cdot y\cdot z$ where $y$ is a product of
    $n(w)$ pairwise commuting generators, whence Parts (1), (2) and (5) of Proposition~\ref{prop:Facts} imply that
    $\la(w)\ge\la(y)=n(w)$.  
\end{proof}

\begin{corollary}
    \label{coro:layerSizeBound}
    Suppose that $\la(w)=2$ \textup{(}so $w$ is FC by Proposition~\ref{prop:a2impliesFC}\textup{)}. 
    \begin{enumerate}
        \item An antichain in $H(w)$ has at most two elements,  and every antichain in $H(w)$ with two elements is
             maximal.
        \item Every layer in the Cartier--Foata form of $w$ has at most two generators.
        \item We have $\abs{\calr(w)}\le 2$.
    \end{enumerate}
\end{corollary}
\begin{proof} 
    We must have $n(w)\le 2$ by Proposition~\ref{prop:a=n}.(1). This implies (1). By the definition of heaps, each layer in
    the Cartier--Foata form of $w$ corresponds to an antichain that has as many elements as the support of the
    layer, so (1) implies (2). Finally, the set $\calr(w)$ equals the
    support of the first layer of $w$ by Remark~\ref{rmk:heapsAndWords}; therefore (3) follows from (2).
\end{proof}

Antichains of heaps appear frequently in this paper. We highlight one property of maximal antichains now. Recall
that an \emph{ideal} in a poset $P$ is a set $\cali$ such that if $y\in \cali$ and $x<y$ for $x,y\in P$ then $x\in
\cali$ ; dually, a \emph{filter} in $P$ is a set $\calf$ such that if $y\in \calf$ and $x>y$ for $x,y\in P$ then
$x\in \calf$. Any set $A\se P$ \emph{generates} an ideal 
\[ \cali_A:=\{i\in H:i\le j\text{ for some $j\in A$}\} \]
and a filter 
\[ \calf_A:=\{i\in H:i\ge j\text{ for some $j\in A$}\}.\] 
If $A$ is an antichain, then any element in the set $P\setminus (\cali_A\cup\calf_A)$, if it exists, can be
adjoined to $A$ to form a larger antichain. This implies the following:

\begin{lemma}
    \label{lemm:maxAntichain} If $A$ is a maximal antichain in $P$, then $P=\cali_A\cup \calf_A$.
\end{lemma}

\subsection{Star operations} 
\label{sec:StarOperations}

In this subsection we discuss star operations on $W$. We emphasize the so-called right star operations and omit
details on the analogous left star operations in the spirit of Remark~\ref{rmk:symmetry}.  

A right star operation is a function from $W$ to $W$ partially defined with respect to a pair of noncommuting
generators $I=\{s,t\}\se S$. Let $m=m(s,t)$ and consider the parabolic subgroup $W_I=\ip{s,t}$. Then each element
$w$ in $W$ admits a unique reduced factorization $w=w^I \cdot w_I$ where $\calr{(w^I)}\cap I = \emptyset$ and
$w_I\in W_I$.  This factorization is the \emph{left coset decomposition} of $w$ with respect to $I$; see \cite[\S
2.4]{BB}. Depending on the value of $w_I$, one of the following mutually exclusive conditions must hold; the
sequences in~(2) and (3) are both infinite if $m=\infty$.
\begin{enumerate} 
    \item $w_I=1$, or $w_I$ is the longest element $sts\dots$ of length $m(s,t)$ in $W_I$; 
    \item $w$ is one of the $(m-1)$ elements \[ x_1:= w^I\cdot s, \quad x_2:= w^I\cdot st,\quad   x_3:= w^I\cdot
         sts,\quad \dots \] where $\call(w_I)=\{s\}$ and $l(w_I)<m$; 
    \item $w$ is one of the $(m-1)$ elements \[ y_1:= w^I\cdot t,\quad y_2:= w^I\cdot ts,\quad  y_3:= w^I\cdot tst,
         \quad\dots \] where $\call(w_I)=\{t\}$ and $l(w_I)<m$.
\end{enumerate} 

\noindent
By definition, we may apply a \emph{right upper star operation with respect to $I$} on $w$ if and only if $w$ is an
element of the form $x_i$ or $y_i$ for some $1\le i\le m-2$; in these cases, the result of the operation is denoted
by $w^*$ and defined to be $x_{i+1}$ or $y_{i+1}$, respectively. Similarly, a \emph{right lower star operation with
respect to $I$} is defined on $w$ precisely when $w$ is an element of the form $x_i$ or $y_i$ for some $2\le i\le
m-1$, whence the result of the operation is the element $w_*:=x_{i-1}$ or $w_*:=y_{i-1}$, respectively.  Left star
operations are defined similarly, and we denote the result of applying a left upper or lower star operation on $w$
by ${}^* w$ or ${}_* w$, respectively. 

When $m=m(s,t)=3$, the (generalized) star operations defined above recover the original star operations introduced
by Kazhdan and Lusztig in \cite{KL}. In this case at most one of $w^*$ and $w_*$ is defined depending on the value
of $w_I$; therefore we may unambiguously speak of \emph{the right star operation with respect to $I$}, without
indicating if the operation is upper or lower. More precisely, both $w^*$ and $w_*$ are undefined if $w_I\in
\{1,sts=tst\}$, $w^*$ is defined while $w_*$ is not if $w_I\in\{s,t\}$, and $w_*$ is defined while $w^*$ is not if
$w_I\in \{st,ts\}$. In the last two cases, we denote whichever one of $w^*$ and $w_*$ is
defined by $w*$,
placing the star sign in the middle. We call the operation $w\mapsto w*$ a \emph{simple star operation}; it is an
involution in the sense that $(w*)*=w$ whenever $w*$ is defined. Similarly, the pair $I=\{s,t\}$ gives rise to a
partially defined involution $w\mapsto *w$ on $W$ that we call the \emph{simple left star operation with respect to
$I$}.

\begin{example}
    \label{eg:StarOperation} 
    Let $(W,S)$ be a Coxeter system with $S=\{a,b,c\}$, $m(a,b)=3$, $m(b,c)=4$ and $m(a,c)=2$. Let $I=\{a,b\},
    J=\{b,c\}$, and let $w=abcab$. Then with respect to $I$, the coset decompositions of $w$ are given by \[ w=
    w_I\cdot {}^I w= aba\cdot cb, \quad w= w^I \cdot w_I = abc\cdot ab, \] so $w_*=abc\cdot a$ while $w^*,
    {}_* w, {}^* w$ are not defined.  Moreover, since $m(a,b)=3$, we have $w*=w_*=abca$ while $*w$ is undefined.
    With respect to $J$, we have \[ w=w_J\cdot {}^J w=b\cdot abcb, \quad w=w^J\cdot w_J=ba\cdot bcb,
    \] so
    ${}^{*}w=cb\cdot abcb, w_*=ba\cdot bc,$ while ${}_*w$ and $w^*$ are not defined.  
\end{example}

Star operations are intimately related to \kl cells. First, left and right cells are closed under left and right
star operations, respectively: 
\begin{prop}[{\cite[Proposition 3.3]{GreenXu}}]
    \label{prop:StarOpAndCells} 
    Let $(W,S)$ be a Coxeter system. Then the following holds with respect to every pair of noncommuting generators
    $\{s,t\}$ in  $S$:
    \begin{enumerate}
        \item  If $w\in W$ and ${}_* w$ is defined, then $w\sim_L {}_* w$;
        \item  If $w\in W$ and $w_*$ is defined, then $w\sim_R w_*$.
    \end{enumerate}
\end{prop}

\noindent Second, simple star operations preserve cell equivalence in the following way: 
\begin{prop}[{\cite[Corollary 4.3]{KL}}]
    \label{prop:SimpleStarOpAndCells} 
    Let $(W,S)$ be a Coxeter system, let $y,w\in W$, and suppose $m(s,t)=3$ for some generators $s,t\in S$.  Then
    the following hold with respect to the pair $\{s,t\}$:
    \begin{enumerate}
        \item if  $y\sim_L w$ and $y*,w*$ are defined, then $y*\sim_L w*$;
        \item if $y\sim_R w$ and $*y,*w$ are defined, then $*y\sim_R *w$. 
    \end{enumerate}
\end{prop}

\noindent 
Later, we will often combine Propositions~\ref{prop:StarOpAndCells} and~\ref{prop:SimpleStarOpAndCells} with
Proposition~\ref{prop:Facts} to show two elements are in the same cell or in distinct cells. Note that
Proposition~\ref{prop:StarOpAndCells} implies that star operations preserve $\la$-values:

\begin{corollary}
    Let $x,y\in W$. If $y$ is obtained from $x$ via any star operation, then $\la(y)=\la(x)$.
    \label{coro:starA}
\end{corollary}
\begin{proof}
    This is immediate from Proposition~\ref{prop:StarOpAndCells} and Proposition~\ref{prop:Facts}.(5).
\end{proof}

\noindent
Our discussion of star operations so far applies to arbitrary elements of $W$, rather than only FC elements. On the
other hand, if an element $w\in W$ is FC, then the heap of $w$ provides a simple characterization of what star
operations can be applied to $w$. We explain why this is true for right lower star operations below. Let
$I=\{s,t\}$ be a pair of noncommuting generators in $S$ as before. By definition, the element $w_*$ with respect to
$I$ is defined and given by $w_*=ws$ if and only if in the coset decomposition $w=w^I\cdot w_I$ we have $2 \le
l(w_I)< m(s,t)$ and the reduced word of $w_I$ is of the form  $w=\dots ts$. The last condition holds if and only if
$s \in\calr(w)$ and $t\in \calr(ws)$, in which case $w$ has a reduced word ending in $s$ and removal of that $s$
results in a reduced word of $w_*$. Remark~\ref{rmk:heapsAndWords} now implies the following result:
\begin{prop}
    \label{prop:Reducibility} Let $w\in W$ be FC and suppose $3\le m(s,t)$ for some $s,t\in S$.  Then $w$ admits a
          right lower star operation with respect to $\{s,t\}$ which results in $ws$ if and only if the conditions
          below all hold:
    \begin{enumerate}
        \item the heap $H(w)$ contains a maximal element $i$ labeled by $s$;
        \item the element $i$ covers an element $j\in H(w)$ labeled by $t$;
        \item the element $i$ is the unique maximal element covering $j$ in $H(w)$.
    \end{enumerate}
\end{prop}

\noindent 
Note that Condition (3) is equivalent to the condition that $j$ becomes maximal upon removal of $i$ from $H(w)$.
The proposition allows us to tell what lower star operations applies to an FC element from a glance at its heap.
For example, given the element $w=abcadb$ and the embedding of $H(w)$ considered in Figure~\ref{fig:HeapExample},
we note that the two highest vertices in the figure are the only maximal elements in $H(w)$ and that only the
removal of top vertex labeled by $b$ results in a new maximal vertex, namely, the higher of the two vertices
labeled by $a$. It follows that only one right lower star operation is defined on $w$, namely, the operation with
respect to $\{a,b\}$ that removes the right descent $b\in \calr(w)$. 

\section{Kazhdan--Lusztig Cells via Stubs}
\label{KLcells}

Recall the following definition of {stubs} from the introduction:
\begin{definition} 
\label{def:stubs} Let $(W,S)$ be a Coxeter system and let $w\in W$. We call $w$ a \emph{left stub} if no right
      lower star operation can be applied to $w$. Similarly, we call $w$ a \emph{right stub} if $w$ admits no left
      lower star operation. We denote the set of left stubs of $\la$-value 2 in $W$ by $\sw$, and denote the set of
      right stubs of $\la$-value 2 in $W$ by $\cals'(W)$.
\end{definition} 

In this section, we assume that $(W,S)$ is an irreducible $\la(2)$-finite Coxeter system, describe
the stubs in $\sw$, and use stubs to find the left, right and two-sided cells in $W_2$. Our main
results on cells are Theorems~\ref{thm:stubDescription}, ~\ref{thm:cellParameterization},
\ref{thm:CellViaDecomp} and~\ref{thm:twoSidedCells}. As we study 2-cells in \autoref{sec:2cells}, we
also introduce two equivalence relations on $\sw$ that will play a crucial role as we study 0-cells
in \autoref{sec:0cells}.  

We note that while some results in the section apply to arbitrary Coxeter systems, stubs seem to
control cells in $W_2$ in the nicest ways only for $\la$(2)-finite systems; see
Remark~\ref{rmk:a2finiteAssumption}.  We also note that since $W_2$ certainly contains no cell when it is
empty, we shall often assume that $(W,S)$ is \emph{nontrivially $\la(2)$-finite} in the sense that
$W_2$ is finite but nonempty.  By Proposition~\ref{prop:facts01}.(4) and the results of \cite{GreenXu}, the
nontrivially $\la(2)$-finite systems are exactly those of types $A_n$ $(n\ge 3)$, $B_n$ $(n\ge 3)$, $\tilde
C_{n-1}$ $(n\ge 5)$, $E_{q,r}$ $(r\ge q\ge 1)$, $F_n$ $(n\ge 4)$ and $H_n$ $(n\ge 3)$. Our description of stubs
and cells will therefore often be given type by type. 

We will often use the prefix ``$\la(2)$-'' as an adjective to indicate an object has $\la$-value 2.
For example, the set $\sw$ is exactly the set of left $\la(2)$-stubs in $W$, and an $\la(2)$-cell
means a cell of $\la$-value 2.

\subsection{Description of stubs}
\label{sec:StubDescriptions} 
We study {$\la(2)$-stubs} in this subsection. As $\la(2)$-elements are FC, we start with the
following observation:

\begin{prop}
    \label{prop:stubCharacterization}
    Let $(W,S)$ be an arbitrary Coxeter system and let $w$ be an FC element in $W$. Let $w=w_p\dots w_2w_1$ be the
    Cartier--Foata form of $w$. Then $w$ is a left stub if and only if every generator $t\in \supp(w_2)$ fails to
    commute with at least two generators $s,s'\in \supp(w_1)$. 
\end{prop}

\begin{remark}
    Note that the above condition on $t$ is vacuously true when $w=w_1$, i.e. when $w$ is a product of commuting
    generators.
\end{remark}

\begin{proof}[Proof of Proposition~\ref{prop:stubCharacterization}]
    Let $t\in \supp(w_2)$ and $s\in \supp(w_1)$. Then $s\neq t$ by
    Proposition~\ref{prop:FCCriterion}.(1), so $t$ is covered by $s$ in $H(w)$ if and only if $m(s,t)\ge 3$.  The
    proposition then follows from Remark~\ref{rmk:heapsAndWords} and Proposition~\ref{prop:Reducibility}.
\end{proof}
\noindent 
The proposition shows that whether an FC element is a left stub is determined ``locally'', by only the first two
layers of its Cartier--Foata form. 

\begin{corollary}
    \label{coro:firstTwoLayers}
    Let $(W,S)$ be an arbitrary Coxeter system.  Let $G$ be the Coxeter diagram of $(W,S)$.  Then an
    $\la(2)$-element $w\in W$ is a left stub if and only if in the Cartier--Foata form $w=w_p\dots w_2w_1$ of $w$,
    we have 
    \begin{enumerate} 
        \item  $w_1=ss'$ for two commuting generators $s, s'$ in $S$;
        \item Every generator $t\in \supp(w_2)$ is adjacent to both $s$ and $s'$ in $G$.  
    \end{enumerate}
\end{corollary}

\begin{proof}
    By Proposition~\ref{prop:stubCharacterization}, it suffices to show that $\supp(w_1)$ contains exactly two
    generators.  By Corollary~\ref{coro:layerSizeBound}, it further suffices to show that
    $\supp(w_1)$ is not empty or a singleton.  This follows from Proposition~\ref{prop:facts01}: if $w_1$ is
    empty, then $w=1$ and $\la(w)=\la(1_W)=0$; if $\supp(w_1)$ contains only one generator $s$, then
    $w_2$ must be empty by Proposition~\ref{prop:stubCharacterization}, which forces $w=w_1=s$ and
    $\la(w)=\la(s)=1$.    
\end{proof}

For nontrivially $\la(2)$-finite systems, Corollary~\ref{coro:firstTwoLayers} turns out to impose strong
restrictions on the first two layers of $\la(2)$-stubs.  The following observation imposes further
restrictions on the deeper layers. 

\begin{lemma}
    \label{lemm:deeperLayers}
    Let $(W,S)$ be an arbitrary Coxeter system. Let $G$ be the Coxeter diagram
    of $(W,S)$. Let $w\in \fc(W)$, let $w=w_p\dots w_1$ be the Cartier--Foata form of $w$, and let~$2<i\le p$. 
    \begin{enumerate}
        \item Every element $t\in \supp(w_i)$ is adjacent in $G$ to some element $s\in \supp(w_{i-1})$.  
        \item If $s,t\in S$ are generators connected by a simple edge in $G$ such that $s$ appears in $w_{i-2}$ and
             $w_{i-1}=t$, then $s$ does not appear in $w_i$.  
    \end{enumerate}
\end{lemma}

\begin{proof}
    Part~(1) simply paraphrases Condition (2) from the characterization of the Cartier--Foata form
    in \autoref{sec:FC}. To prove (2), note that if $s$ appears in $w_{i}$ then we may commute the
    $s$ in $w_{i-2}$ to the left and the $s$ in $w_i$ to the right, if necessary, until they are
    both next to the $t$ in $w_{i-1}$. This results in a reduced factorization $w=x \cdot sts\cdot
    y$ of $w$, contradicting the word criterion for full commutativity. 
\end{proof}

The following example shows in detail how Conditions (1)--(2) from Corollary~\ref{coro:firstTwoLayers} and
Lemma~\ref{lemm:deeperLayers} strongly restrict the Cartier--Foata forms of left $\la(2)$-stubs in a
nontrivially $\la(2)$-finite Coxeter system.  Roughly speaking, the restrictions are strong because
the Coxeter diagram of such a system always contains few pairs of vertices that share a neighbor as
well as few heavy edges.

\begin{example}
    \label{eg:layerAnalysis} 
    We find all left $\la(2)$-stubs in the Coxeter system $(W,S)$ of type $B_n$ in this example.
    Table~\ref{tab:B4stubs} contains the heaps of the six stubs in $\cals(B_4)$, which will serve as an example. 

    Let $w\in\sw$ and let $w=w_p\dots w_2w_1$ be the Cartier--Foata form of $w$.  By Condition (1)
    of Corollary~\ref{coro:firstTwoLayers}, we must have either $w_1=ij$ for some $1\le i,j\le n$ where
    $j>i+2$ or $w_1=i(i+2)$ for some $1\le i\in n-1$. In the former case---which holds for only the
    stub in the third column in Table~\ref{tab:B4stubs} if $n=4$---the second layer of $w$ cannot exist
    by Condition (2) of Corollary~\ref{coro:firstTwoLayers} since $i,j$ share no neighbor in $G$.  In the
    latter case---which corresponds to the first two columns of Table~\ref{tab:B4stubs}---either $w_2$
    does not exist and $w=w_1=i(i+2)$ (first row), or $w_2$ exists and is given by $w_2=(i+1)$, the
    only common neighbor of $i$ and $(i+2)$ in $G$ (second row). In the latter subcase, if $i>1$
    then we have $m(i+1,i)=m(i+1,i+2)=3$, hence Lemma~\ref{lemm:deeperLayers} rules out the possibility
    of a third layer in $w$ and forces $w=w_2w_1=(i+1)\cdot i(i+2)$; if $i=1$, then
    Lemma~\ref{lemm:deeperLayers} implies that we have either $w=w_{2}w_1=2\cdot (13)$ or
    $w=w_3w_2w_1=1\cdot 2\cdot (13)$. In particular, if $w$ contains at least three layers then we
    have to have $w_1=13, w_2=2$ and $w_3=1$, whence $w$ cannot have a fourth layer $w_4$: an
    element $t\in \supp(w_2)$ has to be adjacent to the only generator $1$ in $w_3$ and thus has to
    be $t=2$, so if $w$ contains a fourth layer then $w=\dots\cdot 2\cdot 1\cdot 2\cdot
    (13)=\dots\cdot (2121)\cdot 3$, contradicting the fact that $w$ is FC. The stub $1\cdot 2\cdot
    (13)$ is shown in the third row and first column in Table~\ref{tab:B4stubs}.

    In summary, we can find all left $\la(2)$-stubs in $B_n$, and among them there is a unique stub
    with more than two layers, namely $1\cdot 2\cdot (13)$. The problem of finding left
    $\la(2)$-stubs in the Coxeter system $A_n$ is similar but easier.
\end{example}

\begin{table}
    \begin{tabular}{|c|c|c|}
        \hline
        \begin{tikzpicture}
            \node[main node] (1) {};
            \node[main node] (2) [right=1cm of 1] {};
            \node (11) [above = 0.1cm of 1] {\tiny{$1$}};
            \node (22) [above = 0.1cm of 2] {\tiny{$3$}};
        \end{tikzpicture}
        &
        \begin{tikzpicture}
            \node[main node] (1) {};
            \node[main node] (2) [right=1cm of 1] {};
            \node (11) [above = 0.1cm of 1] {\tiny{$2$}};
            \node (22) [above = 0.1cm of 2] {\tiny{$4$}};
        \end{tikzpicture} 
        &
        \begin{tikzpicture}
            \node[main node] (1) {};
            \node[main node] (2) [right=1cm of 1] {};
            \node (11) [above = 0.1cm of 1] {\tiny{$1$}};
            \node (22) [above = 0.1cm of 2] {\tiny{$4$}};
        \end{tikzpicture}
        \\
        \hline
        \begin{tikzpicture}
            \node[main node] (1) {};
            \node[main node] (2) [right=1cm of 1] {};
            \node[main node] (3) [below right = 0.5cm and 0.5cm of 1] {};
            \node (11) [above = 0.1cm of 1] {\tiny{$1$}};
            \node (22) [above = 0.1cm of 2] {\tiny{$3$}};
            \node (33) [below = 0.1cm of 3] {\tiny{$2$}};
            \path[draw]
            (1)--(3)--(2); 
        \end{tikzpicture}
        &
        \begin{tikzpicture}
            \node[main node] (1) {};
            \node[main node] (2) [right=1cm of 1] {};
            \node[main node] (3) [below right = 0.5cm and 0.5cm of 1] {};
            \node (11) [above = 0.1cm of 1] {\tiny{$2$}};
            \node (22) [above = 0.1cm of 2] {\tiny{$4$}};
            \node (33) [below = 0.1cm of 3] {\tiny{$3$}};
            \path[draw]
            (1)--(3)--(2);  
        \end{tikzpicture}
        &        
        \multicolumn{1}{c}{}\\
        \cline{1-2}
        \begin{tikzpicture} 
            \node[main node] (1) {};
            \node[main node] (2) [right=1cm of 1] {};
            \node[main node] (3) [below right = 0.5cm and 0.5cm of 1] {};
            \node[main node] (4) [below = 1cm of 1] {};
            \node (11) [above = 0.1cm of 1] {\tiny{$1$}};
            \node (22) [above = 0.1cm of 2] {\tiny{$3$}};
            \node (33) [below = 0.1cm of 3] {\tiny{$2$}};
            \node (11) [below = 0.1cm of 4] {\tiny{$1$}};
            \path[draw]
            (1)--(3)--(2)
            (3)--(4);
        \end{tikzpicture}            
        &
        \multicolumn{2}{c}{}\\
        \cline{1-1}
    \end{tabular}
    \vspace{1em} 
    \caption{The six stubs of $B_4$}
    \label{tab:B4stubs}
\end{table}

We classify the left $\la(2)$-stubs of all nontrivially $\la(2)$-finite Coxeter systems in the
following theorem.  In the theorem, we express every element in its Cartier--Foata form and use
$\cdot$ to separate the different layers, and the symbol $\sqcup$ denotes disjoint union. We will
draw the heaps of typical stubs immediately after the proof of the theorem.  The classification lays
the foundation for the subsequent study of $\la(2)$-cells. 

\begin{thm}
    \label{thm:stubDescription}
    Let $(W,S)$ be an irreducible and nontrivially $\la(2)$-finite Coxeter system. Suppose that $(W,S)$ is of type
    $X$, and let $G$ be the Coxeter diagram of $(W,S)$, labeled in the convention explained in Remark~\ref{rmk:labels}.
    Let $\sx$ be the set of all left stubs of $\la$-value 2 in $W$.  Then $\sx$ can be described as follows. 
    \begin{enumerate}
        \item If $X=A_n$ for some $n\ge 3$, then $\sx=\mathcal{S}_1\sqcup\mathcal{S}_2$ where 
            \[
            \mathcal{S}_1=\{x_{ij}:=ij\,\vert\, 1\le i,j\le n,j>i+1\}\]
            and
            \[
                \mathcal{S}_2=\{y_i:=i\cdot
                (i-1)(i+1)\,\vert\, 1< i<n\}.
            \]
        \item If $X=B_n$ for some $n\ge 3$, then $\sx=\mathcal{S}_1\sqcup\mathcal{S}_2\sqcup \cals_3$ where
            $\mathcal{S}_1,\mathcal{S}_2$ are the same \textup{(}see Remark
           ~\ref{rmk:stubRemarks1}.\textup{(}1\textup{)}\textup{)} as in \textup{(}1\textup{)}
            and 
            \[
                \cals_3=\{z_1:=1\cdot 2\cdot 13\}.
            \] 
        \item If $X=\tilde C_{n-1}$ for some $n\ge 5$, then $\sx =\cals_1\sqcup \cals_2\sqcup
             \cals_3\sqcup\cals_4$ where $\cals_1,\cals_2,\cals_3$ are the same as in
             \textup{(}2\textup{)} and 
            \[
                \cals_4=\{z_n:=n\cdot (n-1)\cdot (n-2)n\}.
            \]
        \item If $X=E_{q,r}$ for some $r\ge q\ge 1$, then $\sx=\cals_1\sqcup\cals_2\sqcup
             \mathcal{T}_2\sqcup\cals_3$ for the sets 
            \[
            \cals_1=\{x_{st}\,\vert\,s,t\in S, m(s,t)=2\},\]
            \[
                \cals_2=\{y_{i}:= i\cdot (i-1)(i+1)\,\vert\, -q< i<r\},
            \]
            \[
                \mathcal{T}_2=\{y'_0:=0\cdot (-1)v\}\sqcup\{y''_0:=0\cdot 1v\},
            \]
            and $\cals_3=\{z_s: s\in S, s\neq 0\}$ where 
            \[
                z_s= 
                \begin{cases}
                    s\cdot (s-1)\cdot \ldots \cdot 0\cdot (-1)v &\text{if $s\neq v$ and $1\le s\le
                    r$},\\
                    s\cdot (s+1)\cdot \ldots \cdot 0\cdot (1)v & \text{if $s\neq v$ and $-q\le s\le -1$},\\
                    v\cdot 0 \cdot (-1)1&\text{if $s=v$}.
                \end{cases}
            \]
        \item If $X=F_n$ for some $n\ge 4$, then $\sx=\cals_1\sqcup\cals_2\sqcup\cals_3$
            where  $\cals_1,\cals_2$ are the same as in \textup{(}1\textup{)} and 
            \[
                \cals_3=\{
                    z_1:=1\cdot 2\cdot 3\cdot 24, z_2:=2\cdot 3\cdot
                24\}\sqcup\{z_i\,\vert\, 3\le i\le n\}, 
            \]
            where $z_i=i\cdot (i-1)\cdot \ldots\cdot 3\cdot 2\cdot 13$ for each $3\le i\le n$.
        \item If $X=H_n$ for some $n\ge 3$, then $\sx=\cals_1\sqcup\cals_2\sqcup\cals_3\sqcup\cals_4$ where
            $\cals_1,\cals_2,\cals_3$ are the same as in \textup{(}2\textup{)} and 
            \[
                \cals_4=\{z_i:=i\cdot (i-1)\cdot \ldots\cdot 1\cdot 2\cdot 13\,\vert\,
                2\le i\le n\}.
            \]
        \item In each of the above, we listed each stub of $\la$-value 2 exactly once. {\textup{(}}For example, two
            $x_{ij},x_{i'j'}$ from $ \cals_1$ are equal only if $i=i',j=j'$.{\textup{)}} The cardinality of $\sx$ 
            in $W$ is given by 
            \[ 
                \abs{\sx}=
                \begin{cases}\binom{n}{2} -1 & \text{if $X=A_n$ \textup{(}$n\ge 3$\textup{)}};\\ 
                    \binom{n}{2} & \text{if
                      $X=B_n$ \textup{(}$n\ge 3$\textup{)}};\\ 
                      \binom{n}{2} +1 & \text{if $X=\tilde C_{n-1}$
                      \textup{(}$n\ge 5$\textup{)}};\\ 
                      \binom{n+1}{2}-1 & \text{if $X= E_{q,r}$ \textup{(}$q,r\ge 1, n=q+r+2=\abs{S}$\textup{)}};\\ 
                      \binom{n+1}{2}-1 & \text{if $X= F_n$ \textup{(}$n\ge 4$\textup{)}};\\
                      \binom{n+1}{2}-1 & \text{if $X= H_n$ \textup{(}$n\ge 3$\textup{)}}.
                \end{cases} 
            \]
    \end{enumerate}
\end{thm}


\begin{remark}
    \label{rmk:stubRemarks1}
    \begin{enumerate}
        \item In Theorem~\ref{thm:stubDescription}, by saying that a certain set $\cals_i$ in some part  is the same as
             the set $\cals_i$ given from an earlier part, we mean that the former set consists of the elements
             given by the same reduced words as those of the elements in the latter set. 
        \item By symmetry, an element $w\in W$ is right stub if and only if $w\inverse$ is a left stub, so
             $\cals'(W)=\{w\inverse: w\in \cals(W)\}$ for each Coxeter system $(W,S)$ appearing in Theorem
            ~\ref{thm:stubDescription}.
    \end{enumerate} 
\end{remark}


\begin{proof}[Proof of Theorem~\ref{thm:stubDescription}]
    Part~(7) follows from Parts (1)--(6) by inspection (to see that every element in $\sw$ is indeed
    listed only once) and by easy counting arguments, so it suffices to show that in each of Cases
    (1)--(6), an element $w\in W$ is a left \twostub if and only if $w\in \sx$. The ``only if''
    implication follows from analysis of the layers of left $\la(2)$-stubs:
    Corollary~\ref{coro:firstTwoLayers} implies that the first two layers of such a stub $w$ must satisfy
    Conditions (1)--(2) given in the corollary, whence Lemma~\ref{lemm:deeperLayers} forces $w$ to be an
    element in $\sw$ by arguments similar to those we used for type $B_n$ in
    Example~\ref{eg:layerAnalysis}. The ``if'' implication, i.e. the fact that every element $w\in \sw$ is
    a left $\la(2)$-stub, can be proven as follows: first, draw the heap of the specified reduced
    word of $w$ we gave (such as $x_{13}=13\in A_n$) and use Proposition~\ref{prop:FCCriterion} to see that $w$
    is indeed FC; second, use Proposition~\ref{prop:Reducibility} to see that $w$ is a stub;
    finally, note that we can transform $w$ to its first layer $w_1$ by a sequence of left lower
    star operations by the analog of Proposition~\ref{prop:Reducibility} for left star operations,
    and then invoke Corollary~\ref{coro:starA} to conclude that $\la(w)=\la(w_1)=2$.  The proof is complete.
\end{proof}

Let us describe and draw the heaps of stubs described in Theorem~\ref{thm:stubDescription}. Call a stub
$w\in \sw$ \emph{short} if $l(w)=2$, \emph{medium} if $l(w)=3$, and \emph{long} if $l(w)>3$. Then
the short, medium, and long stubs in $\sw$ are precisely those labeled by ``$x$'', ``$y$''
(including $y'_0$ and $y''_0$ in type $E_{q,r}$) and ``$z$'' in Theorem~\ref{thm:stubDescription}; they are
also precisely those with one, two, and more than two layers, respectively.  The heaps of the short
and medium stubs take the forms shown in Table~\ref{tab:xyStubs}.  Long stubs exist in all
$\la(2)$-finite Coxeter systems except $A_n$.  In type $B,\tilde C, F$ and $H$, every long stub
involves a ``turn'' at the second layer, that is, there is always a heavy edge $\{s,t\}$ in the
Coxeter diagram such that $s$ appears in the first and third layers of the stub while $t$ appears in
the second layer. The stub $1213\in B_4$ pictured in Table~\ref{tab:B4stubs} is such an example.
Table~\ref{tab:zBCFH} contains more typical such stubs.  Finally, in type $E_{q,r}$, each long stub $w$
with Cartier--Foata form $w=w_p\dots w_2w_1$ satisfies the conditions $\supp(w_2)=\{0\}$,
$\supp(w_1)=\{s,t\}$, $\supp(w_2)=\{u\}$ where $\{s,t,u\}=\{-1,1,v\}$; the heaps of the long stubs are
shown in Table~\ref{tab:zE}.

    \begin{table}
        \begin{tabularx}{0.8\textwidth}{|X|X|}
            \hline
            \raisebox{-1.2em}{
    \begin{tikzpicture}
        \node (0) {short stubs:};
        \node [main node] (1) [right = 0.3cm of 0] {};
        \node [main node] (2) [right = 1cm of 1] {};
\end{tikzpicture}
}
    &
    \raisebox{-2em}{
    \begin{tikzpicture}
        \node (00) {medium stubs:};
        \node [main node] (4) [above right = 0.05cm and 0.3cm of 00] {};
        \node [main node] (5) [below right = 0.5cm and 0.5cm of 4] {};
        \node [main node] (6) [above right = 0.5cm and 0.5cm of 5] {};
        \node (x) [below =1.5em of 6] {};
        \path[draw]
        (4)--(5)--(6);
    \end{tikzpicture}
}
    \\
    \hline
\end{tabularx}
\vspace{1em} 
\caption{The short and medium stubs of $\sw$}
    \label{tab:xyStubs}
    \end{table}


    \begin{table}
        \centering
        \begin{tabular}{|c|c|}
            \hline
            \raisebox{0.8em}{
                \begin{tikzpicture}
                    \node (0) {$z_{n}\in \tilde C_{n-1}$:};
                    \node [main node] (1) [right = 0.8cm of 0] {};
                    \node [main node] (2) [above left = 0.5cm and 0.5cm of 1] {};
                    \node [main node] (3) [above right = 0.5cm and 0.5cm of 1] {};
                    \node [main node] (4) [below right = 0.5cm and 0.5cm of 1] {};

                    \node (11) [right = 0.01cm of 1] {\tiny{$(n-1)$}};
                    \node (22) [above = 0.01cm of 2] {\tiny{$(n-2)$}};
                    \node (33) [above = 0.01cm of 3] {\tiny{$n$}};
                    \node (44) [below = 0.01cm of 4] {\tiny{$n$}};
                    \path[draw]
                    (2)--(1)--(3)
                    (1)--(4);
            \end{tikzpicture}}
            & 
            \raisebox{-0.1em}{
                \begin{tikzpicture}
                    \node (x0) {$z_{1}\in F_{n}$:};
                    \node [main node] (x1) [right = 1.3cm of x0] {};
                    \node [main node] (x2) [above left = 0.5cm and 0.5cm of x1] {};
                    \node [main node] (x3) [above right = 0.5cm and 0.5cm of x1] {};
                    \node [main node] (x4) [below left = 0.5cm and 0.5cm of x1] {};
                    \node [main node] (x5) [below left = 0.5cm and 0.5cm of x4] {};

                    \node (x11) [below = 0.01cm of x1] {\tiny{$3$}};
                    \node (x22) [above = 0.01cm of x2] {\tiny{$2$}};
                    \node (x33) [above = 0.01cm of x3] {\tiny{$4$}};
                    \node (x44) [below = 0.01cm of x4] {\tiny{$2$}};
                    \node (x55) [below = 0.01cm of x5] {\tiny{$1$}};
                    \path[draw]
                    (x2)--(x1)--(x3)
                    (x1)--(x4)--(x5);
                \end{tikzpicture} 
            }
            \\ \hline

            \raisebox{0.2em}{
                \begin{tikzpicture}
                    \node (y0) {$z_i\in F_n$:};
                    \node (b) [below = 0.01cm of y0] {\tiny{$(3\le i\le n)$}};
                    \node [main node] (y1) [above right = 0.1cm and 0.8cm of y0] {};
                    \node [main node] (y2) [above left = 0.5cm and 0.5cm of y1] {};
                    \node [main node] (y3) [above right = 0.5cm and 0.5cm of y1] {};
                    \node [main node] (y4) [below right = 0.5cm and 0.5cm of y1] {};
                    \node [main node] (y5) [below right = 0.8cm and 0.8cm of y4] {};
                    \node [main node] (y6) [below right = 0.5cm and 0.5cm of y5] {};

                    \node (y11) [right = 0.01cm of y1] {\tiny{$2$}};
                    \node (y22) [above = 0.01cm of y2] {\tiny{$1$}};
                    \node (y33) [above = 0.01cm of y3] {\tiny{$3$}};
                    \node (y44) [right = 0.01cm of y4] {\tiny{$3$}};
                    \node (y44) [right = 0.01cm of y5] {\tiny{$(i-1)$}};
                    \node (y44) [right = 0.01cm of y6] {\tiny{$i$}};
                    \path[draw]
                    (y2)--(y1)--(y3)
                    (y1)--(y4)
                    (y5)--(y6);
                    \path[draw,dashed]
                    (y4)--(y5);
            \end{tikzpicture}}
            & 
            \begin{tikzpicture}
                \node (yx0) {$z_{i}\in H_{n}$:};
                \node (b) [below = 0.01cm of yx0] {\tiny{$(1\le i\le n)$}};
                \node [main node] (yx1) [right =  0.8cm of yx0] {};
                \node [main node] (yx2) [above left = 0.5cm and 0.5cm of yx1] {};
                \node [main node] (yx3) [above right = 0.5cm and 0.5cm of yx1] {};
                \node [main node] (yx4) [below left = 0.5cm and 0.5cm of yx1] {};
                \node [main node] (yx5) [below right = 0.5cm and 0.5cm of yx4] {};
                \node [main node] (yx6) [below right = 0.8cm and 0.8cm of yx5] {};
                \node [main node] (yx7) [below right = 0.5cm and 0.5cm of yx6] {};

                \node (yx11) [right = 0.01cm of yx1] {\tiny{$2$}};
                \node (yx22) [above = 0.01cm of yx2] {\tiny{$1$}};
                \node (yx33) [above = 0.01cm of yx3] {\tiny{$3$}};
                \node (yx44) [right = 0.01cm of yx4] {\tiny{$1$}};
                \node (yx55) [right = 0.01cm of yx5] {\tiny{$2$}};
                \node (yx66) [right = 0.01cm of yx6] {\tiny{$(i-1)$}};
                \node (yx77) [right = 0.01cm of yx7] {\tiny{$i$}};
                \path[draw]
                (yx2)--(yx1)--(yx3)
                (yx1)--(yx4)--(yx5)
                (yx6)--(yx7);
                \path[draw,dashed]
                (yx5)--(yx6);
            \end{tikzpicture}\\ \hline
        \end{tabular}
        \vspace{1em} 
        \caption{Typical long stubs of $\sw$ in $B_n, \tilde C_{n-1}, F_n$ and $H_n$}
        \label{tab:zBCFH}
    \end{table}

    \begin{table}
        \begin{tabularx}{\textwidth}{|X|c|X|}
            \hline
            \begin{tikzpicture}
                \node (y0) {$z_i$:};
                \node (b) [below = 0.01cm of y0] {\tiny{$(i\in \Z_{<0})$}};
                \node [main node] (y1) [right = 1.8cm of y0] {};
                \node [main node] (y2) [above left = 0.5cm and 0.5cm of y1] {};
                \node [main node] (y3) [above right = 0.5cm and 0.5cm of y1] {};
                \node [main node] (y4) [below left = 0.5cm and 0.5cm of y1] {};
                \node [main node] (y5) [below left = 0.8cm and 0.8cm of y4] {};
                \node [main node] (y6) [below left = 0.5cm and 0.5cm of y5] {};

                \node (y11) [right = 0.01cm of y1] {\tiny{$0$}};
                \node (y22) [above = 0.01cm of y2] {\tiny{$1$}};
                \node (y33) [above = 0.01cm of y3] {\tiny{$v$}};
                \node (y44) [right = 0.01cm of y4] {\tiny{$-1$}};
                \node (y44) [right = 0.01cm of y5] {\tiny{$(i+1)$}};
                \node (y44) [right = 0.01cm of y6] {\tiny{$i$}};
                \path[draw]
                (y2)--(y1)--(y3)
                (y1)--(y4)
                (y5)--(y6);
                \path[draw,dashed]
                (y4)--(y5);
            \end{tikzpicture}
            &
            \raisebox{1.5em}{
                \begin{tikzpicture}
                    \node (0) {$z_{v}$:};
                    \node [main node] (1) [right = 0.8cm of 0] {};
                    \node [main node] (2) [above left = 0.5cm and 0.5cm of 1] {};
                    \node [main node] (3) [above right = 0.5cm and 0.5cm of 1] {};
                    \node [main node] (4) [below  = 0.5cm  of 1] {};

                    \node (11) [right = 0.01cm of 1] {\tiny{$0$}};
                    \node (22) [above = 0.005cm of 2] {\tiny{$-1$}};
                    \node (33) [above = 0.01cm of 3] {\tiny{$1$}};
                    \node (44) [below = 0.01cm of 4] {\tiny{$v$}};
                    \path[draw]
                    (2)--(1)--(3)
                    (1)--(4);
                \end{tikzpicture}
            }
            & 
            \begin{tikzpicture}
                \node (y0) {$z_i$:};
                \node (b) [below = 0.01cm of y0] {\tiny{$(i\in \Z_{>0})$}};
                \node [main node] (y1) [right = 0.8cm of y0] {};
                \node [main node] (y2) [above left = 0.5cm and 0.5cm of y1] {};
                \node [main node] (y3) [above right = 0.5cm and 0.5cm of y1] {};
                \node [main node] (y4) [below right = 0.5cm and 0.5cm of y1] {};
                \node [main node] (y5) [below right = 0.8cm and 0.8cm of y4] {};
                \node [main node] (y6) [below right = 0.5cm and 0.5cm of y5] {};

                \node (y11) [right = 0.01cm of y1] {\tiny{$0$}};
                \node (y22) [above = 0.005cm of y2] {\tiny{$-1$}};
                \node (y33) [above = 0.01cm of y3] {\tiny{$v$}};
                \node (y44) [right = 0.01cm of y4] {\tiny{$1$}};
                \node (y44) [right = 0.01cm of y5] {\tiny{$(i-1)$}};
                \node (y44) [right = 0.01cm of y6] {\tiny{$i$}};
                \path[draw]
                (y2)--(y1)--(y3)
                (y1)--(y4)
                (y5)--(y6);
                \path[draw,dashed]
                (y4)--(y5);
            \end{tikzpicture}\\\hline
        \end{tabularx}
        \vspace{1em} 
        \caption{The long stubs of $\cals(E_{q,r})$}
        \label{tab:zE}
    \end{table}



We record a few features of the stubs in Theorem~\ref{thm:stubDescription} for future use:

\begin{lemma}
    \label{lemm:stubRemarks}
    Let $(W,S)$ be an arbitrary $\la(2)$-finite Coxeter group from Theorem~\ref{thm:stubDescription} and keep the notation
    of the theorem.  Then the stubs in $\sw$ have the following properties.  
    \begin{enumerate}
        \item The left descent set of each short stub  $x_{tt'}$, medium stub $y_t$ \tul including
             $y'_0,y''_0$ in type $E_{q,r}$\tur and long stub $z_t$ equals $\{t,t'\}, \{t\}$ and
             $\{t\}$, respectively. In particular, a stub $w\in \sw$ has two left descents if and
             only if it is a short stub, and in this case we can recover $w$ from $\call(w)$: if
             $\call(w)=\{t,t'\}$ then $w=x_{tt'}$.
        \item Let $w\in \sw$ be a medium or long stub and let $w=w_p\dots w_1$ be its Cartier--Foata
             factorization.  Then $l(w_1)=2$ and  $l(w_j)=1$ for all $2\le j\le p$, so we may pick a
             reduced word $ \ul{w}= s_1\dots s_{k}(s_{k+1}s_{k+2}) $ of $w$ such that
             $w_1=s_{k+1}s_{k+2}$, $w_2=s_{k}, w_3=s_{k-1}$, etc.
        \item Let $w$ and $\ul w$ be as in \textup{(2)}. Then in the heap $H(\ul w)$, every element
             $i<k+1$ is comparable both to $k+1$ and to $k+2$ via the convex chains of coverings
             $i\prec i+1\prec \dots k\prec k+1$ and $i\prec i+1\prec \dots k\prec k+2$. In
             particular, the element $k$ is both the only element covered by $k+1$ and the only
             element covered by $k+2$ in $H(\ul w)$. 
        \item Let $w$ and $\ul w$ be as in \textup{(2)}. Then we may start with $w$ and successively
             apply left lower star operations with respect to the noncommuting pairs of generators
             $\{s_1,s_2\}, \{s_2,s_3\},\dots, \{s_{k},s_{k+1}\}$ to obtain the sequence $w,
             s_2s_3\dots s_k(s_{k+1}s_{k+2}), s_3\dots s_{k}(s_{k+1}s_{k+2}), \dots,$
             $s_k(s_{k+1}s_{k+2}), w_1$. All elements in this sequence lie in $\sw$ and share the
             same left cell (and hence the same right descents) as $w$. In particular, $w$ is
             in the same left cell and has the same right descents as its first layer $w_1 =
             s_{k+1}s_{k+2}$, which is itself a short stub.
            \end{enumerate}
    \end{lemma}
    \begin{proof} All the claims are readily confirmed by inspection of Theorem~\ref{thm:stubDescription}
          and the pictures in Tables~\ref{tab:xyStubs}--\ref{tab:zE}.     
    \end{proof}

    \begin{remark}
      \label{rmk:a.vs.n}
      The question of whether the $n$-value (as defined in the paragraph before Proposition~\ref{prop:a=n}) and
      $\la$-value of an FC element always agree in a general Coxeter system is open. The results of
      this subsection allow us to enhance Proposition~\ref{prop:a=n}.(1) and give a partial answer to this
      question in the following way: 

\begin{prop} 
\label{prop:a.vs.n}
Let $(W,S)$ be an arbitrary Coxeter system and let $w\in \fc(W)$.  \begin{enumerate} \item We
   have $\la(w)=0$ if and only if $n(w)=0$.  \item We have $\la(w)=1$ if and only if
   $n(w)=1$.
\item If $(W,S)$ is $\la(2)$-finite, then $\la(w)=2$ if and only if $n(w)=2$.  \end{enumerate}
  \end{prop} 
\begin{proof} (1) By Proposition~\ref{prop:facts01}.(1), the equations $\la(w)=0$ and $n(w)=0$ both hold
     exactly when $w$ is the identity element and has an empty reduced word, so they are
     equivalent.

(2) If $\la(w)=1$, then $n(w)\neq 0$ by Part~(1) and $n(w)\le 1$ by Proposition~\ref{prop:a=n}.(1), and
therefore $n(w)=1$.  Conversely, if $n(w)=1$ then the heap $H(w)$ is a chain, so no reduced word
of $w$ contains two consecutive generators that commute. The definition of FC elements then
implies that $w$ has a unique reduced word, so $\la(w)=1$ by Proposition~\ref{prop:facts01}.(2). It follows
that $\la(w)=1$ if and only if $n(w)=1$.

(3) Suppose $(W,S)$ is $\la(2)$-finite. If $\la(w)=2$, then $n(w)\notin\{0,1\}$ by Parts (1) and~(2)
and $n(w)\le 2$ by Proposition~\ref{prop:a=n}.(1), and therefore $n(w)=2$. Now suppose that~$n(w)=2$. We prove that
$\la(w)=2$ by induction on the length $l(w)$ of $w$. We have $l(w)= \abs{H(w)}\ge n(w)=2$, so in the
base case we have $l(w)=2$. Since $n(w)=2$, we must have $w=st$ for two commuting generators $s,t\in
S$ in this case; thus $\la(w)=2$ by Proposition~\ref{prop:Facts}.(2), as desired.  If $l(w)>2$, then $w$
either admits or does not admit a right lower star operation. In the former case, if a right lower
star operation takes $w$ to some element $w'$ then we have $\la(w)=\la(w')=n(w')=n(w)=2$, where the
first equality holds by Corollary~\ref{coro:starA}, the second equality holds by induction, and the third
equality holds because star operations do not change $n$-values of FC elements (see the proof of
Proposition 3.15 in \cite{GreenXu}). In the latter case, the element $w$ is a left stub, so
Corollary~\ref{coro:firstTwoLayers} implies that the first two layers of $w$ satisfy Conditions (1)--(2) from
the corollary. As pointed out in the proof of Theorem~\ref{thm:stubDescription}, this forces $w$ to be an
element in the set $\sw$ given in the theorem, and therefore $w$ indeed has $\la$-value~2.  It
follows by induction that $\la(w)=2$. We have proved that $\la(w)=2$ if and only if~$n(w)=2$.
\end{proof}
    \end{remark}

\subsection{Stubs parameterize 1-cells}
\label{sec:1cells}
Let $(W,S)$ be a \ntatf Coxeter system throughout this subsection. We characterize 1-cells in $W_2$
via stubs in two ways, first in terms of star operations in Theorem~\ref{thm:cellParameterization} and then
in terms of reduced words and the weak Bruhat order in Theorem~\ref{thm:CellViaDecomp}. To start, we show
that distinct left stubs lie in distinct right cells in the following two results.

\begin{lemma}
    \label{lemm:stubDescents}
    Let $u,w$ be two distinct left $\la(2)$-stubs in $W$.  Then we either have $\call(u)\neq \call(w)$ or have
    $\call(u)=\call(w)=\{s\}$ for some $s\in S$. Moreover, in the latter case the elements $u':=su$ and $w':=sw$
    are also stubs, and we have $\call(u')\neq \call(w')$. 
\end{lemma}
\begin{proof} 
    We use Lemma~\ref{lemm:stubRemarks}. Suppose that $\call(u)=\call(w)$.  Since $u\neq w$ by assumption and left stubs
    with two left descents can be recovered from the descent by Part~(1) of the lemma, we must have
    $\call(u)=\call(w)=\{s\}$ for some $s\in S$. This proves the first claim.  For the second claim, note that in
    the notation of Theorem~\ref{thm:stubDescription}, we can have $\call(u)=\call(w)=\{s\}$ for distinct stubs $u,w$ only
    in the following cases where $(W,S)$ is of type $E_{q,r}, F_n$ or $H_n$; outside these types, the claim holds
    vacuously.

    \begin{enumerate}[label=\textup{(\roman*)}]
        \item \textbf{Case 1:} $(W,S)$ is of type $E_{q,r}$, and we have either one of two subcases:
            \begin{enumerate}                
                \item $s=0$ and $u,w$ are two $y$-stubs from the set $\{y_0,y'_0,y''_0\}$ where \[y_0=0\cdot
                    x_{(-1)1}, \quad y_0'=0\cdot x_{(-1)v},\quad y''_0=0\cdot x_{1v};\]
                \item $s=i$ for some nonzero integer $i$ with $-q<i<r$, and 
                    \[ 
                        \{u,w\}=\{y_i=i\cdot
                    x_{(i-1)(i+1)},\; z_i=i\cdot z'\} \] 
                    where $z'$ is a long stub.
            \end{enumerate}
        \item\textbf{Case 2:} 
            $(W,S)$ is of type $F_n$, and there is some $1<i<n$ such that 
            \[ \{u,w\}=\{y_i=i\cdot x_{(i-1)(i+1)},\; z_i=i\cdot z'\} \]
            where $z'$ is a long stub.
        \item\textbf{Case 3:} $(W,S)$ is of type $H_n$, and there is some $1<i<n$ such that 
            \[ \{u,w\}=\{y_i=i\cdot x_{(i-1)(i+1)},\; z_i=i\cdot z'\}  \]
            where $z'$ is a long stub.
    \end{enumerate} 
   We have $\call(u')\neq \call(w')$ for the stubs  $u'=su$ and $w'=sw$ in Case 1(a) by inspection,
   and the same is true in all the other cases by Lemma~\ref{lemm:stubRemarks}.(1) because one of $u',w'$
   is a short stub while the other is not.
\end{proof}

\begin{prop}
    \label{prop:distinctCells}
    Let $u,w$ be distinct left $\la(2)$-stubs of $W$. Then $u\not\sim_R w$. 
\end{prop}
\begin{proof}
    Keep the notation from Theorem~\ref{thm:stubDescription}, Lemma~\ref{lemm:stubDescents}, and the proof of
    the lemma.  We have $u\not\sim_R w$ if $\call(u)\neq \call(w)$ by Proposition~\ref{prop:Facts}.(3), so
    Lemma~\ref{lemm:stubDescents} implies that it suffices to treat Cases 1--3 from its proof, where
    $\call(u)=\call(w)=\{s\}$ for some $s\in S$. We may show $u\not\sim_R w$ by finding a generator
    $t\in S$ such that $m(s,t)=3$ and $\call(*u)\neq \call(*w)$, where $*$ denotes the simple left
    star operation with respect to~$\{s,t\}$: the fact that $\call(*u)\neq \call(*w)$ implies that
    $*u\not\sim_R *w$, so $u\not\sim_R w$ by Proposition~\ref{prop:SimpleStarOpAndCells}.(2). We explain how to
    find such a generator $t$ below.

    In Case 1(a) of the proof of Lemma~\ref{lemm:stubDescents}, we may take $t$ to be the unique
    generator appearing in both the stubs $u'=su$ and $w'=sw$.  For example, if $u=y_0$ and
    $w=y_0'$, then $u'= x_{(-1)1}, w'=x_{(-1)v}$ and we may take $t=-1$, whence $*u=u', *w=w'$ and
    $\call(*u)\neq \call(*w)$.  In Cases 1(b), 2 and 3, we note that $s$ has a numerical label $i$
    and that $i-1,i+1$ are also generators in $S$. By our labeling of generators
    (Remark~\ref{rmk:labels}), it follows that $m(i,i+1)=3$ except when $(W,S)$ is of Coxeter type $F_n$,
    $i=2$ and $\{u,w\}=\{y_2=2(13),z_2=2\cdot 3\cdot (24)\}$.  In this exceptional case we pick
    $t=i-1=1$, so that $m(s,t)=3$ and $\{*u,*w\}=\{13=x_{13},1\cdot 2\cdot 3\cdot (24)=z_1\}$ with
    respect to~$\{s,t\}$. In all other cases we may pick $t=i+1$, whence $m(s,t)=3$ and
    $\{*u,*w\}=\{*y_i,*z_i\}=\{x_{(i-1)(i+1)},z_{i+1}\}$ with respect to~$\{s,t\}$. Here, the fact
    that $*z_i=z_{i+1}$ in the last set can be easily seen from heaps in Tables~\ref{tab:zBCFH} and
   ~\ref{tab:zE}, with the star operation on $z_i$ being a lower one if and only if $(W,S)$ is of
    type $E_{q,r}$ and $-q<i<0$. In all cases, one of $*u$ and $*w$ is an short stub and the other
    is a long stub, so $\call(*u)\neq \call(*w)$ by Lemma~\ref{lemm:stubRemarks}.(1), as desired.
\end{proof}

We are ready to describe cells in terms of stubs and star operations.

\begin{definition}
    \label{def:closure}
    For each element $w\in W$, we define the \emph{right upper star closure} of $w$ to be the set of
    all elements $y\in W$ for which there exists a sequence $z_1=w, z_2,\dots, z_q=y$ such that
    $z_{i+1}$ can be obtained from $z_{i}$ via a right upper star operation for all $1\le i\le q-1$;
    we denote the set by $R_w$. Similarly, for each $w\in W$ we define its \emph{left upper star
    closure} to be the set $L_w$ containing all elements that can be obtained from $w$ via a
    sequence of left upper star operations.
\end{definition}

\begin{thm}
    \label{thm:cellParameterization}
    Let $W$ be a nontrivially $\la(2)$-finite Coxeter group, let $W_2$ be the set of $w\in W$ with $\la(w)=2$, and let
    $\cals(W)$ be the set of left $\la(2)$-stubs in $W$.  Then the set $R_w$ forms a right \kl cell
    for every $w\in \sx$.  Moreover, we have  $R_w\cap R_{w'}=\emptyset$ for distinct
    $w,w'\in\cals(W)$, and $W_2=\sqcup_{w\in \cals(W)}R_w$. In particular, the number of right cells
    in $W_2$ equals the cardinality of $\sw$ given in Part \textup{(}7\textup{)} of
    Theorem~\ref{thm:stubDescription}.
\end{thm}

\begin{proof}
    Let $R(w)$ be the right \kl cell containing $w$ for each $w\in W$. Then we have $R_w\se R(w)$ by
    Proposition~\ref{prop:StarOpAndCells}.  On the other hand, by the definition of stubs we have
    $W_2=\cup_{w\in \sw} R_w$.  It follows that to prove the theorem it suffices to show that
    $w\not\sim_R w'$ whenever $w,w'$ are distinct elements in $\sw$. This holds by
    Proposition~\ref{prop:distinctCells}.
\end{proof}

Using the fact that the set $R_w$ is a right cell for each $w\in\sw$, we now work towards a second
description of the cells, in terms of reduced words.  


\begin{definition}
    \label{def:stubDecomp}
    Let $w\in W$. We define a \emph{left stub decomposition} of $w$ to be a reduced factorization of
    the form $w=x\cdot z$ where $x$ is a left stub and $z\in W$.  Similarly, we define a \emph{right
    stub decomposition} of $w$ to be a reduced factorization of the form $w=z\cdot y$ where $y$ is a
    right stub.  
\end{definition}

\begin{thm}
    \label{thm:CellViaDecomp}
    Let $x\in \sw$ and let $w$ be an element of $\la$-value 2 in $W$. Let $R(x)$ be the right cell containing $x$
    \textup{(}so that $R(x)=R_x$ by Theorem~\ref{thm:cellParameterization}\textup{)}.
    \begin{enumerate}
        \item We have  $w\in R_x$ if and only if $w$ has a left stub decomposition of the form $w=x\cdot z$.
        \item We have $w\in R(x)$ if and only if $x\le^R w$, i.e. we have 
        \[
            R(x)=\{z\in W_2: x\le^R z\}.
        \]
        \item The element $w$ has a unique left stub decomposition in the sense that the stub $x'$ and the
             element  $z'$ in the decomposition $w=x'\cdot z'$ are both unique.
    \end{enumerate}
\end{thm}

\noindent
Note that Part~(2) of the theorem, which characterizes the right cell in $W_2$ in terms of the right
weak Bruhat order, may be viewed as an analog of Proposition~\ref{prop:facts01}.(3). To prove the theorem, we
will use the following lemma. 

\begin{lemma}
    \label{lemm:decompAndDescent}
    Let $(W,S)$ be an arbitrary Coxeter system and let $w\in \fc(W)$.
    \begin{enumerate}
        \item Suppose that $w$ has a reduced word of the form 
            \[
                \ul w=s_1\dots s_k (s_{k+1}s_{k+2}\dots s_{k+j})s_{k+j+1}\dots s_{q}
         \]
         where the set $A:=\{k+1,k+2,\dots,k+j\}$ forms a maximal antichain in
             the heap $H(w)$. Let $x=s_1\dots s_k(s_{k+1}\dots s_{k+j})$ and $y=(s_{k+1}\dots s_{k+j})s_{k+j+1}\dots s_q$.
             Then we have $H(x)=\cali_A$ and $H(y)=\calf_A$ as sets. 
            
        \item In the above setting, we have $\call(w)=\call(x)$ and $\calr(w)=\calr(y)$.  

        \item Let $x\in \sw$ and let $w$ be an element of $\la$-value 2 with a
    left stub decomposition of the form $w=x\cdot z$. Then $\call(w)=\call(x)$.
    \end{enumerate}
\end{lemma}
\begin{proof}
    (1) Since $A$ is a maximal antichain in $H(w)$, we have $H(\ul w)=\cali_{A}\cup\calf_{A}$ by
    Lemma~\ref{lemm:maxAntichain}. For any $k+j+1\le i\le q$, we have $i\notin\cali_A$ by the definition
    of heaps; therefore $i\in \calf_A$. Similarly, we have $i\in \cali_A$ for all $k+1\le i\le q$.
    It follows that $H(x)=\cali_A$ and $H(y)=\calf_A$.

    (2) Part~(1) implies that the set of minimal elements of $H(x)$ and $H(w)$ coincide, so
    $\call(w)= \call(x)$ by Remark~\ref{rmk:heapsAndWords}. Similarly we have $\calr(w)=\calr(y)$.

    (3) By Lemma~\ref{lemm:stubRemarks}.(2), there is a reduced word $\ul x=s_1\dots
    s_{k}(s_{k+1}s_{k+2})$ of $x$ where $(s_{k+1}s_{k+2})$ equals the first layer of $x$. Let
    $s_{k+3}\dots s_{q}$ be a reduced word of $z$, so that $\ul{w}:=s_1\dots s_k\cdot
    (s_{k+1}s_{k+2})\cdot s_{k+3}\dots s_{q}$ is a reduced word of $w$.  The set $A:=\{k+1,k+2\}$
    forms a maximal antichain in the heap $H(\ul w)$ by Corollary~\ref{coro:layerSizeBound}.(1), and
    therefore $\call(w)=\call(x)$ by Parts (1) and (2).
\end{proof}

\begin{proof}[Proof of Theorem~\ref{thm:CellViaDecomp}]
    Part~(1) implies Part~(2) by the definition of $\le^R$ and Proposition~\ref{prop:Facts}.(5). Part~(3) also
    follows from Part~(1): if $w=x''\cdot z''$ is another left stub decomposition of $w$ then $w\in
    R(x')$ and $w\in R(x'')$, and therefore $x\sim_R w\sim_R x''$; this forces $x'=x''$, and hence
    $z'=z''$, by Proposition~\ref{prop:distinctCells}.

    It remains to prove Part~(1). The ``only if'' implication follows from the definition of $R_x$
    and upper star operations. To prove the ``if'' implication, suppose that $w$ has a stub
    factorization of the form $w=x\cdot z$, and let $y\in \sw$ be the unique stub such that $w\in
    R_y$. The unique existence of such a stub is guaranteed by Theorem~\ref{thm:cellParameterization}, and
    we need to prove that $x=y$.

    Since $w\in R_y$, we have $w\sim_R y$ and hence  $\call(y)=\call(w)$ by Proposition~\ref{prop:Facts}.(3).
    On the other hand, since $w=x\cdot z$, we have $\call(x)=\call(w)$ by
    Lemma~\ref{lemm:decompAndDescent}. It follows that $\call(x)=\call(y)$, where $1\le
    \abs{\call(x)}=\abs{\call(y)}\le 2$ by the analog of Corollary~\ref{coro:layerSizeBound}.(3) for left
    descents. If $\abs{\call(x)}=\abs{\call(y)}=2$, then $x=y$ by Lemma~\ref{lemm:stubRemarks}.(1).  If
    $\abs{\call(x)}=\abs{\call(y)}=1$, say with $\call(x)=\call(y)=\{s\}$, then $l(x),l(y)>2$ and
    the elements $x':=sx$ and $y'=sy$ are stubs with $\call(x')\neq \call(y')$ by
    Lemma~\ref{lemm:stubRemarks}.(4). The element $w':=sw$ has stub decompositions of the form
    $w'=x'\cdot z'$ and $w'=y'\cdot z''$, which forces $\call(x')=\call(w')=\call(y')$ by
    Lemma~\ref{lemm:decompAndDescent}. Since $\call(x)=\call(y)$ and $\call(x')=\call(y')$,
    Lemma~\ref{lemm:stubDescents} implies that $x=y$.  The proof is complete.  
\end{proof}

